% Sumber LaTeX untuk buku ``Metode Aljabar'' dalam bahasa Mandarin
% Hak Cipta 2018  Wen-Wei Li.
% Izin diberikan untuk menyalin, mendistribusikan, dan/atau memodifikasi
% dokumen ini berdasarkan ketentuan Creative Commons
% Attribution 4.0 International (CC BY 4.0)
% http://creativecommons.org/licenses/by/4.0/

% Untuk disertakan
\chapter{Ekstensi Medan}\label{sec:field-ext}
Bab ini membahas struktur umum medan, tetapi tidak menyentuh grup Galois. Pokok dasar yang dikaji dalam teori medan ialah pembenaman medan, atau dengan kata lain, ekstensi medan. Perhatian utama kita tertuju pada ekstensi aljabar, yang secara alami berkaitan dengan sifat-sifat polinomial dan akar-akarnya. Sudut pandang yang dipakai buku ini ialah memanfaatkan keberadaan penutup aljabar secara sistematis untuk menangani ekstensi aljabar seluas mungkin. Karena ekstensi medan $F$ kita pandang sebagai suatu jenis khusus $F$-aljabar, konsep-konsep seperti polinomial minimal, teras, dan norma dalam suatu ekstensi telah dibahas pada Bab Tujuh; bab ini hanya memberikan ulasan yang diperlukan tanpa mengulangi pembuktiannya.

\begin{wenxintishi}
	Perangkat dasar untuk menangani ekstensi aljabar ialah keberadaan pembenaman medan dan perluasan pembenaman. Pada akhirnya, semua pernyataan ini menyangkut sifat-sifat polinomial dan akar-akarnya, sehingga teori medan sama sekali tidak abstrak. Pada pembacaan pertama, \S\S\ref{sec:pins}--\ref{sec:tensor-field} dapat dilewati. Pembahasan lengkap mengenai medan berhingga ditunda hingga bab berikutnya.
\end{wenxintishi}

\section{Beberapa Jenis Ekstensi}
Medan adalah gelanggang pembagian komutatif. Suatu subgelanggang dari medan yang dengan sendirinya juga merupakan medan disebut submedan. Karena medan tidak memiliki ideal sejati nontrivial, setiap homomorfisme gelanggang $F \to E$ di antara medan $E$ dan $F$ bersifat injektif; dengan kata lain, pelaku utama teori medan ialah pembenaman medan $F \hookrightarrow E$. Pembenaman yang sama dapat dipandang dari berbagai sudut:
\begin{compactitem}
	\item memberikan pembenaman $u: F \to E$ setara dengan memberikan struktur $F$-aljabar pada $E$;
	\item jika $F$ diidentifikasi dengan $u(F) \subset E$, maka $F$ dapat dipandang sebagai submedan dari medan $E$;
	\item dalam situasi tersebut, $E$ juga disebut suatu \emph{ekstensi} atau \emph{medan ekstensi} dari $F$. Buku ini menggunakan lambang $E|F$, sedangkan banyak literatur menggunakan $E/F$.\index{yukuozhang@ekstensi medan (field extension)}\index[sym1]{E/F@$E \mid F$, $E/F$}
\end{compactitem}
Perhatikan bahwa melalui perkalian, $E$ secara alami menjadi ruang vektor atas $F$. Buku ini tidak mensyaratkan bahwa dalam ekstensi medan $E|F$, medan $F$ merupakan submedan dari $E$ dan $F \hookrightarrow E$ merupakan peta inklusi, meskipun anggapan ini sangat memudahkan dan juga sah dalam pengertian isomorfisme berikut.
 
\begin{convention}\index{yukuozhang!pembenaman (embedding)}
	Untuk ekstensi medan $E|F$ dan $E'|F$, suatu pembenaman-$F$ di antara keduanya didefinisikan sebagai homomorfisme gelanggang $\phi: E \to E'$ yang membuat diagram berikut komutatif:
	\[\begin{tikzcd}[column sep=small, row sep=small]
		E \arrow[rr, "\phi"] & & E' \\
		& F \arrow[hookrightarrow, lu] \arrow[hookrightarrow, ru] &
	\end{tikzcd}\]
	Semua $F$-pembenaman $E \to E'$ membentuk himpunan $\Hom_F(E, E')$; ini tidak lain ialah himpunan-$\Hom$ dalam kategori $F$-aljabar. Dengan pola yang sama, dapat didefinisikan konsep isomorfisme dan automorfisme di antara medan-medan ekstensi dari $F$.
\end{convention}
Untuk ekstensi $E|F$ seperti di atas, jika dari $E$ suatu subaljabar atas $F$ dengan sendirinya merupakan medan, maka subaljabar itu disebut \emph{subekstensi} dari $E|F$. Dari sudut pandang pembenaman, subekstensi dari $u: F \hookrightarrow E$ ialah submedan di dalam $E$ yang memenuhi $E' \supset u(F)$; karena itu, bila $F \subset E$, subekstensi dari $E|F$ juga disebut \emph{medan antara}.\index{zhongjianyu@medan antara (intermediate field)}

Untuk medan $F$, karakteristiknya $\text{char}(F)$ didefinisikan sebagaimana dalam Definisi \ref{def:ring-characteristic}: $\text{char}(F)$ membangkitkan $\Ker[\Z \to F]$. Karakteristik medan adalah nol atau suatu bilangan prima $p$. Sesuai dengan itu, submedan terkecil dari medan $F$ ialah $\Q$ (dalam karakteristik nol) atau $\F_p := \Z/p\Z$ (dalam karakteristik prima $p$), dan disebut \emph{submedan prima} dari $F$. Setiap submedan dari $F$ merupakan subekstensi dari $F$ atas submedan primanya. Untuk setiap $u: F \hookrightarrow E$ berlaku $\text{char}(E)=\text{char}(F)$, dan pembatasan $u$ pada submedan prima merupakan peta identitas. \index{tezheng}\index[sym1]{char}\index{suziyu}

\begin{definition}
	Misalkan $E|F$ suatu ekstensi medan.
	\begin{itemize}
		\item Derajat $E|F$ didefinisikan sebagai kardinal $[E:F] := \dim_F E$;\index{yukuozhang!derajat (degree)}\index[sym1]{$[E:F]$}
		\item jika $[E:F]$ berhingga, maka $E|F$ disebut \emph{ekstensi berhingga};\index{yukuozhang!berhingga}
		\item jika setiap $x \in E$ merupakan unsur aljabar atas $F$ (Definisi \ref{def:algebraic-element}), yaitu jika terdapat $P \in F[X] \smallsetminus \{0\}$ sehingga $P(x)=0$, maka $E|F$ disebut \emph{ekstensi aljabar}.\index{yukuozhang!aljabar}
		\item Untuk $x \in E$, definisikan $F[x]$ sebagai $F$-subaljabar yang dibangkitkan oleh $x$ (lihat pembahasan pada awal \S\ref{sec:integrality-finiteness}); kasus beberapa unsur, $F[x,y, \ldots]$, didefinisikan dengan cara serupa.
		\item Selanjutnya, definisikan $F(x)|F$ sebagai subekstensi yang dibangkitkan oleh $x$ atas $F$. Unsur-unsur medan $F(x)$ berbentuk $\frac{P(x)}{Q(x)}$, dengan $P,Q \in F[X]$ dan $Q(x) \neq 0$; kasus beberapa unsur, $F(x,y,\ldots)$, didefinisikan dengan cara serupa.\index[sym1]{$F(x,y,\ldots)$}
		\item Selanjutnya, jika suatu keluarga unsur $\{x_i\}_{i \in I}$ di dalam $E$ memenuhi $F(x_i : i \in I) = E$, maka $\{x_i\}_{i \in I}$ disebut himpunan pembangkit dari $E|F$. Ekstensi yang memiliki himpunan pembangkit berhingga disebut \emph{ekstensi yang dibangkitkan secara berhingga}, sedangkan yang dapat dibangkitkan oleh satu unsur disebut \emph{ekstensi sederhana}.\index{yukuozhang!sederhana}
	\end{itemize}
\end{definition}

Ekstensi berhingga $E|F$ tentu dibangkitkan secara berhingga: jika $x_1, \ldots, x_k$ merupakan basis $E$ sebagai ruang vektor atas $F$, maka jelas $E=F(x_1, \ldots, x_k)$. Selain itu, $[E:F]=1 \iff F \rightiso F \cdot 1 = E$, dan hal ini sering langsung ditulis $E=F$.

\begin{example}\label{eg:field-basic-examples}
	Sebelum mendalami teori umum, kita tinjau dua contoh klasik yang berasal dari abad ke-19.\index{shuyu@medan bilangan (number field)}\index{hanshuyu@medan fungsi (function field)}
	\begin{itemize}
		\item Menurut definisi, bilangan aljabar ialah unsur $\CC$ yang memenuhi persamaan polinomial $P(\alpha)=0$, dengan $P \in \Q[t]$ tak nol ($t$ menyatakan peubah). Teori bilangan mempelajari medan ekstensi yang diperoleh dengan menambahkan pada $\Q$ sejumlah berhingga bilangan aljabar $\alpha_1, \ldots, \alpha_n$, yakni $\Q(\alpha_1, \ldots, \alpha_n)|\Q$ dalam notasi kita, serta sifat-sifat seperti pembenaman, derajat, dan automorfisme di antara medan-medan tersebut. Segera akan kita buktikan bahwa semua ekstensi ini merupakan ekstensi aljabar berhingga, dan bahwa seluruh bilangan aljabar membentuk, di dalam $\CC$, submedan $\overline{\Q}$. Operasi penambahan unsur $F \leadsto F(\alpha, \beta, \ldots)$ semacam ini merupakan salah satu cara kerja khas teori medan.
		\item Misalkan $X$ suatu manifold kompleks mulus, kompak, dan berdimensi kompleks satu; sebagai manifold riil, objek ini berdimensi dua dan juga disebut permukaan Riemann kompak. Semua fungsi meromorfik pada $X$ membentuk medan fungsi $\CC(X)$ dari $X$. Dapat dibuktikan bahwa suatu peta holomorfik surjektif $f: X \to Y$ di antara permukaan Riemann kompak menginduksi pembenaman medan $f^*: \CC(Y) \hookrightarrow \CC(X)$ melalui tarikan balik fungsi meromorfik. Derajat ekstensi medan $\CC(X)|\CC(Y)$ tepat sama dengan derajat peta $f$ dalam arti geometrik. Bila $Y$ adalah garis projektif kompleks $\mathbb{P}^1$, maka $\CC(Y)$ diidentifikasi dengan medan fungsi rasional satu peubah $\CC(t)$, sedangkan ekstensi berhingganya $\CC(X)$ disebut medan fungsi aljabar kompleks. Sebuah teorema dasar dalam teori kurva aljabar menyatakan bahwa pembenaman-$\CC$ di antara medan fungsi mencerminkan dengan tepat semua peta holomorfik nontrivial di antara permukaan Riemann. Dari sini dapat diperoleh gambaran kasar mengenai makna geometrik $\CC(X)$.
	\end{itemize}
	Terdapat analogi yang halus di antara medan bilangan dan medan fungsi, yang berawal dari wawasan mendalam Dedekind dan Weber \cite{DW82}.
\end{example}

\begin{definition}\index{yukuozhang!kompositum (compositum)}
	Misalkan $\Omega|F$ suatu ekstensi medan dan $(E_i|F)_{i \in I}$ suatu keluarga subekstensi di dalamnya. \emph{Kompositum} $\bigvee_{i \in I} E_i$ didefinisikan sebagai medan terkecil di dalam $\Omega$ yang memuat semua $E_i$. Unsur-unsurnya berbentuk pecahan rasional
	\begin{gather}\label{eqn:fields-compositum}
		\frac{P(x_{i_1}, \ldots, x_{i_n})}{Q(x_{i_1}, \ldots, x_{i_n})}, \quad Q(x_{i_1}, \ldots, x_{i_n}) \neq 0
	\end{gather}
	dengan $n \geq 0$, $P, Q \in F[X_1, \ldots, X_n]$, $i_1, \ldots, i_n \in I$, dan $x_{i_k} \in E_{i_k}$. Kompositum sejumlah berhingga subekstensi juga ditulis $E_1 E_2$, dan seterusnya.
\end{definition}
Jika semua subekstensi dari $\Omega|F$ disusun sebagai poset menurut $\subset$, maka kompositum memberikan supremum bagi setiap subhimpunannya, sedangkan infimum diberikan oleh irisan subekstensi $\bigcap_{i \in I} E_i$.

Misalkan $L|E$ dan $E|F$ ekstensi medan. Komposisi pembenaman $F \hookrightarrow E \hookrightarrow L$ memberikan ekstensi medan $L|F$. Struktur semacam ini disebut menara ekstensi medan. Pernyataan berikut merupakan penerapan langsung dari Lemma \ref{prop:free-transitivity} dan Korolari \ref{prop:norm-trace-transitivity-2}. \index{yukuozhang!menara (tower)}
\begin{proposition}[Sifat menara untuk derajat]\label{prop:field-tower-degree}
	Untuk $L$, $E$, dan $F$ seperti di atas,
	\begin{compactitem}
		\item derajat ekstensi medan memenuhi $[L:F]=[L:E][E:F]$;
		\item jika $(x_i)_{i \in I}$ dan $(y_j)_{j \in J}$ masing-masing merupakan basis $L$ atas $E$ dan basis $E$ atas $F$, maka $(x_i y_j)_{(i,j) \in I \times J}$ merupakan basis $L$ atas $F$.
	\end{compactitem}
\end{proposition}
Sifat keterbagian derajat memiliki amat banyak penerapan. Sebagai contoh, bila $[L:F]$ adalah bilangan prima, medan antara $E$ harus sama dengan $L$ atau $F$.

Sekarang kita beralih ke salah satu pokok penting bab ini, yaitu kajian ekstensi aljabar. Sifat-sifat berikut merupakan titik awal bagi semua argumen.
\begin{enumerate}
	\item Untuk unsur aljabar $x \in E$ dalam ekstensi $E|F$, dapat didefinisikan \emph{polinomial minimal} $P_x \in F[X]$ atas $F$ (Definisi \ref{def:algebraic-element}); untuk setiap $Q \in F[X]$ berlaku $Q(x) = 0 \iff P_x \mid Q$.
	\item Selanjutnya, Lemma \ref{prop:minimal-polynomial} menyatakan bahwa $P_x$ selalu tak tereduksi dan $F(x) = F[x]$; kita memiliki isomorfisme $F$-aljabar
		\begin{align*}
			F[X]/(P_x) & \longrightiso F(x) \\
			f + (P_x) & \longmapsto f(x) = \text{ev}_x(f), \\
			X + (P_x) & \longmapsto x
		\end{align*}
		Dengan demikian, polinomial minimal secara intrinsik mencirikan struktur medan $F(x)$, dan dari sini juga terlihat bahwa $[F(x):F] = \deg P_x$.
	\item Untuk ekstensi $E|F$ dan setiap subekstensinya $E'|F$, unsur $x \in E$ tentu juga aljabar atas $E'$, dan polinomial minimal $x$ atas $E'$ harus membagi $P_x$; karena itu $[E'(x):E'] \leq [F(x):F]$.
\end{enumerate}
Dalam \S\ref{sec:splitting-field} kita akan memberikan cara umum untuk membangun ekstensi aljabar.

\begin{example}[Ekstensi kuadratik]
	Misalkan $\text{char}(F) \neq 2$. Setiap ekstensi yang memenuhi $[E:F]=2$ harus berbentuk $E = F(\sqrt{a})$, dengan $a \in F^\times \smallsetminus F^{\times 2}$; di sini menurut konvensi, $\sqrt{a} \in E$ menyatakan salah satu akar kuadrat dari $a$. Memang, $[E:F]=2$ mengakibatkan $x \in E \smallsetminus F \implies E = F(x) \simeq F[X]/(P_x)$. Dengan melengkapkan kuadrat pada polinomial minimal $P_x = X^2 + tX + s$ menjadi $(X + \frac{t}{2})^2 + (s - \frac{t^2}{4})$, lebih lanjut dapat diasumsikan bahwa $P_x = X^2 - a$, dan ini memberikan $x = \sqrt{a}$.
\end{example}

\begin{lemma}\label{prop:finiteness-generation}
	Jika ekstensi medan $E|F$ memenuhi $E=F(y_1, \ldots, y_m)$, dengan setiap $y_i$ merupakan unsur aljabar atas $F$, maka $E|F$ merupakan ekstensi berhingga, dan dalam hal ini $E=F[y_1, \ldots, y_m]$.
\end{lemma}
\begin{proof}
	Dari $F(y_1, \ldots, y_m) = F(y_1, \ldots, y_{m-1})(y_m)$, $1 \leq m \leq n$, Proposisi \ref{prop:field-tower-degree} bersama pengamatan di atas memberikan
	\begin{align*}
		[F(y_1, \ldots, y_n):F] & = \prod_{m=1}^n \left[ F(y_1, \ldots, y_m): F(y_1, \ldots, y_{m-1}) \right] \\
		& = \prod_{m=1}^n \left[ F(y_1, \ldots, y_{m-1})(y_m): F(y_1, \ldots, y_{m-1}) \right] \\
		& \leq \prod_{m=1}^n [F(y_m):F] < \infty,
	\end{align*}
	sehingga $E|F$ merupakan ekstensi berhingga. Secara rekursif terhadap $m$ diperoleh
	\begin{equation*}\begin{aligned}
		E & = F(y_1, \ldots, y_{m-1})(y_m) = F(y_1, \ldots, y_{m-1})[y_m] \\
		& = F[y_1, \ldots, y_{m-1}][y_m] = F[y_1, \ldots, y_m],
	\end{aligned}\end{equation*}
	sebagaimana hendak dibuktikan.
\end{proof}

\begin{lemma}\label{prop:finite-vs-algebraic}
	Ekstensi berhingga tidak lain ialah ekstensi aljabar yang dibangkitkan secara berhingga.
\end{lemma}
\begin{proof}
	Teorema \ref{prop:integrality-finiteness} memastikan bahwa ekstensi berhingga $E|F$ merupakan ekstensi aljabar, dan tentu juga dibangkitkan secara berhingga. Arah sebaliknya merupakan akibat langsung dari Lemma \ref{prop:finiteness-generation}.
\end{proof}

Kita kerap perlu mengkaji perilaku sifat-sifat ekstensi di bawah berbagai operasi. Untuk menghemat penulisan, diperkenalkan istilah dari \cite[V. \S 1]{Lang02} berikut.
\begin{convention}\label{con:dist-extensions}\index{yukuozhang!istimewa (distinguished)}
	Misalkan $\mathcal{E}$ suatu kelas ekstensi medan. Kelas $\mathcal{E}$ disebut \emph{istimewa} jika sifat-sifat berikut terpenuhi:
	\begin{enumerate}[\bfseries {D}.1]
		\item untuk setiap menara ekstensi medan $L|E$ dan $E|F$, berlaku $(E|F \in \mathcal{E}) \wedge (L|E \in \mathcal{E}) \iff L|F \in \mathcal{E}$;
		\item misalkan $L|F$ dan $M|F$ subekstensi dari suatu $\Omega|F$ yang diberikan; maka $L|F \in \mathcal{E} \implies LM|M \in \mathcal{E}$;
		\item untuk $L|F$ dan $M|F$ seperti di atas, $(L|F \in \mathcal{E}) \wedge (M|F \in \mathcal{E}) \implies LM|F \in \mathcal{E}$. Perhatikan bahwa sifat ini sebenarnya merupakan akibat dari dua sifat pertama, dan juga mengakibatkan bahwa $\mathcal{E}$ tertutup terhadap kompositum berhingga.
	\end{enumerate}
	Ketiga keadaan di atas sering dinyatakan dengan diagram medan berikut:
	\[ \begin{tikzcd}
		L \arrow[dash,d] \\ E \arrow[dash,d] \\ F
	\end{tikzcd} \quad \begin{tikzcd}[row sep=small, column sep=tiny]
		& LM \arrow[dash,ld] \arrow[dash,rrdd] & & \\
		L \arrow[dash,rrdd] & & & \\
		& & & M \arrow[dash,ld] \\
		& & F &\\
	\end{tikzcd} \quad \begin{tikzcd}[column sep=small]
		& LM \arrow[dash,ld] \arrow[dash,rd] & \\
		L \arrow[dash,rd] & & M \arrow[dash,ld] \\
		& F &
	\end{tikzcd}\]
	Jika \textbf{D.3} diperkuat dengan mensyaratkan bahwa kompositum setiap keluarga subekstensi dari $\Omega|F$ yang termasuk dalam $\mathcal{E}$ tetap termasuk dalam $\mathcal{E}$, termasuk kompositum tak hingga, maka $\mathcal{E}$ disebut \emph{tertutup terhadap kompositum}.
\end{convention}

Tetapkan suatu kelas ekstensi istimewa $\mathcal{E}$ seperti ini. Di dalam suatu ekstensi $\Omega|F$ yang diberikan, kita dapat mengambil gabungan
\[ \bigcup_{\substack{L|F: \; \Omega|F \;\text{mempunyai subekstensi ini} \\ L|F \in \mathcal{E}}} L; \]
Andaikan $F|F \in \mathcal{E}$ sehingga gabungan tersebut tak kosong. Karena $\mathcal{E}$ tertutup terhadap kompositum berhingga, gabungan $\bigcup_{L|F \in \mathcal{E}} L$ tertutup terhadap penjumlahan, perkalian, dan pengambilan invers unsur tak nol, sehingga merupakan submedan dari $\Omega$. Sesungguhnya, gabungan ini juga sama dengan kompositum semua subekstensi $L|F \in \mathcal{E}$.

Syarat pengambilan gabungan dapat diperlonggar. Untuk $\Omega|F$ yang diberikan, jika suatu keluarga subekstensinya $\mathcal{E}$ membentuk poset terarah ke atas menurut $\subset$ dalam arti Definisi \ref{def:filtrant-poset}, maka gabungannya tetap merupakan subekstensi; bila suatu keluarga subekstensi tertutup terhadap kompositum berhingga, keluarga tersebut otomatis membentuk poset terarah ke atas menurut $\subset$. Berikut sebuah contoh.
\begin{example}
	Untuk suatu ekstensi medan $E|F$, tetapkan $\mathcal{E} = \{ E|F\; \text{subekstensi yang dibangkitkan secara berhingga dari} \}$; maksudnya, unsur-unsur $\mathcal{E}$ ialah subekstensi berbentuk $F(a_1, \ldots, a_m)$, dengan $m \in \Z_{\geq 1}$ dan $a_1, \ldots, a_m \in E$. Karena
	\[ F(a_1, \ldots, a_m) F(b_1, \ldots, b_n) = F(a_1, \ldots, a_m, b_1, \ldots, b_n), \]
	terlihat bahwa $\mathcal{E}$ tertutup terhadap kompositum berhingga, tetapi jelas tidak tertutup terhadap kompositum tak hingga. Gabungan subekstensi-subekstensi dalam $\mathcal{E}$ ialah seluruh $E|F$, dengan alasan sederhana bahwa $x \in F(x)$ untuk setiap $x \in E$.
\end{example}

\begin{lemma}\label{prop:finite-ext-dist}
	Ekstensi berhingga bersifat istimewa.
\end{lemma}
\begin{proof}
	Untuk ekstensi $E|F$ dan $L|E$ yang dipertimbangkan dalam \textbf{D.1}, persamaan $[L:F]=[L:E][E:F]$ dari Proposisi \ref{prop:field-tower-degree} mengakibatkan bahwa $L|F$ berhingga jika dan hanya jika $L|E$ dan $E|F$ keduanya berhingga. Sekarang tinjau \textbf{D.2}: ambil subekstensi $L|F$ dan $M|F$ dari $\Omega|F$ sedemikian sehingga $L|F$ berhingga. Kita dapat menulis $L=F(x_1, \ldots, x_n)$, dengan setiap $x_i$ unsur aljabar. Maka $LM = M(x_1, \ldots, x_n)|M$. Setiap $x_i$ juga aljabar atas $M$, sehingga Lemma \ref{prop:finiteness-generation} langsung menunjukkan bahwa $LM|M$ berhingga. Jadi, ekstensi berhingga bersifat istimewa.
\end{proof}

\begin{proposition}\label{prop:alg-ext-dist}
	Ekstensi aljabar merupakan gabungan subekstensi berhingganya. Ekstensi aljabar bersifat istimewa dan tertutup terhadap kompositum.
\end{proposition}
\begin{proof}
	Setiap ekstensi aljabar $E|F$ merupakan gabungan semua $F(x)|F$, dengan $x$ merentang seluruh unsur $E$, sehingga merupakan gabungan subekstensi berhingga.

	Sekarang kita buktikan keistimewaan ekstensi aljabar. Tinjau ekstensi medan $L|E$ dan $E|F$ serta sifat \textbf{D.1}. Jelas bahwa jika $L|F$ aljabar, maka $L|E$ dan $E|F$ keduanya aljabar. Sebaliknya, andaikan $L|E$ dan $E|F$ keduanya ekstensi aljabar. Untuk $x \in L$, misalkan $P_x \in E[X]$ polinomial minimalnya; koefisien-koefisien $P_x$ membangkitkan suatu subekstensi $E_0|F$ yang dibangkitkan secara berhingga, sehingga Lemma \ref{prop:finite-vs-algebraic} memastikan bahwa $[E_0:F]$ berhingga. Karena $[E_0(x) : E_0] \leq \deg P_x$, juga diperoleh $[F(x):F] \leq [E_0(x):F] = [E_0(x):E_0] [E_0:F]$ berhingga, sehingga $x$ aljabar atas $F$.

	Selanjutnya, kita buktikan bahwa kompositum $LM|F$ dari dua subekstensi aljabar $L|F$ dan $M|F$ dari $\Omega|F$ tetap aljabar (sifat \textbf{D.3}). Bila unsur $x \in LM$ ditulis dalam bentuk \eqref{eqn:fields-compositum}, hanya sejumlah berhingga unsur dari $L$ dan $M$ yang terlibat. Karena subekstensi yang dibangkitkan secara berhingga dari $L$ dan $M$ pasti merupakan ekstensi berhingga (Lemma \ref{prop:finite-vs-algebraic}), dan kompositum dua ekstensi berhingga tetap berhingga (Lemma \ref{prop:finite-ext-dist}), maka $x$ aljabar atas $F$. Sifat ini dapat diperluas ke kompositum suatu keluarga sebarang subekstensi aljabar $E_i|F$, sebab setiap $x \in \bigvee_i E_i$ selalu terletak dalam kompositum sejumlah berhingga $E_i$.

	Jika hanya diasumsikan bahwa $L|F$ merupakan ekstensi aljabar, maka karena $LM|M$ adalah kompositum semua $M(x)|M$ ($x \in L$), dan $[M(x):M] \leq [F(x):F] < \infty$, dua langkah sebelumnya menunjukkan bahwa $LM|M$ merupakan ekstensi aljabar. Ini membuktikan \textbf{D.2}.
\end{proof}

\begin{proposition}\label{prop:alg-ext-cardinality}
	Untuk setiap ekstensi aljabar $E|F$ berlaku
	\begin{compactitem}
		\item $|E|=|F|$, jika $|F|$ tak hingga;
		\item $|E| \leq \aleph_0$, jika $|F|$ berhingga.
	\end{compactitem}
	Di sini $\aleph_0$ menyatakan kardinal tak hingga pertama.
\end{proposition}
\begin{proof}
	Misalkan $\text{Irr}(F)_n$ menyatakan subhimpunan polinomial monik tak tereduksi berderajat $n$ dalam $F[X]$, dan $\text{Irr}(F) := \displaystyle\bigcup_{n \geq 1} \text{Irr}(F)_n$. Definisikan peta $m: E \to \text{Irr}(F)$ yang memetakan $x \in E$ ke polinomial minimalnya $P_x \in F[X]$. Banyak akar sebuah polinomial tidak melebihi derajatnya, sehingga aritmetika kardinal (lihat \S\ref{sec:cardinal-number} untuk perincian) memberikan
	\begin{align*}
		|E| & = \sum_{n \geq 1} \left| m^{-1}(\text{Irr}(F)_n) \right| \leq \sum_{n \geq 1} n \left| \text{Irr}(F)_n \right| \\
		& \leq \sum_{n \geq 1} n |F|^{n-1}.
	\end{align*}
	Bila $|F|$ tak hingga, Korolari \ref{prop:cardinal-max} mengakibatkan $n|F|^{n-1} = |F|$, sehingga $|E| \leq |\Z_{\geq 1}| \cdot |F| = \aleph_0 |F| = |F|$. Bila $|F|$ berhingga, Korolari \ref{prop:cardinal-max} yang sama memberikan $|E| \leq |\Z_{\geq 1}| \cdot |\Z_{\geq 1}| = \aleph_0$.
\end{proof}

\section{Penutup Aljabar}
Tujuan awal teori medan adalah mempelajari akar-akar polinomial. Untuk setiap unsur aljabar $x \in E$ dalam ekstensi medan $E|F$, kita telah mendefinisikan polinomial minimal $P_x \in F[X]$ yang sesuai, yaitu polinomial tak tereduksi monik yang memenuhi $P_x(x)=0$. Sebaliknya, dari medan $F$ yang diberikan dan polinomial tak konstan $P \in F[X]$, bisakah kita membangun ekstensi medan $E|F$ sehingga $P$ memiliki akar di $E$? Jelas hanya perlu mempertimbangkan kasus di mana $P$ tak tereduksi. Ide konstruksinya adalah menambahkan akar ke $F$ secara "formal"; Lemma \ref{prop:minimal-polynomial}, melalui isomorfisme $F[X]/(P_x) \xrightarrow{X \mapsto x} F(x)=F[x]$, menunjukkan jalan yang jelas.

\begin{proposition}[L.\ Kronecker]\label{prop:adjunction-root}
	Misalkan $P \in F[X]$ adalah polinomial tak tereduksi, definisikan $F$-aljabar $E := F[X]/(P)$. Maka $E|F$ adalah ekstensi medan, $[E:F] = \deg P$ dan koset $x := X + (P) \in E$ memenuhi $P(x)=0$.
\end{proposition}
\begin{proof}
	Karena $F[X]$ adalah gelanggang ideal utama dan $P$ tak tereduksi, $(P)$ adalah ideal maksimal. Oleh karena itu $E$ memang merupakan medan. Jelas bahwa $E$ sebagai ruang vektor atas $F$ memiliki basis $\{ x^i : 0 \leq i \leq \deg P \}$. Dengan memandang $P$ sebagai unsur $E[X]$, mudah dihitung dalam $E$ bahwa $P(x) = P(X + (P)) = P(X) + (P) = 0$. Terbukti.
\end{proof}

\begin{example}[A.\ L.\ Cauchy, 1847]\label{eg:Cauchy-C}
	Medan kompleks $\CC$ dapat dibangun sebagai kelas-kelas ekuivalensi polinomial berkoefisien riil sebagai berikut. Polinomial kuadrat $X^2+1 \in \R[X]$ tidak memiliki akar riil, sehingga tak tereduksi. Menurut Proposisi \ref{prop:adjunction-root}, bangun ekstensi kuadrat $\R[X]/(X^2+1)$ dari $\R$; tulis koset $X + (X^2+1)$ sebagai $i$. Struktur $\R[X]/(X^2+1) = \R \oplus \R i$ sebagai $\R$-aljabar sepenuhnya ditentukan oleh relasi $i^2 + 1=0$, yang tepat mencirikan medan kompleks. Jadi $\R[X]/(X^2+1) \simeq \CC$.
\end{example}

\begin{proposition}\label{prop:field-embedding}
	Misalkan $E|F$, $L|F$ adalah ekstensi medan.
	\begin{compactenum}[(i)]
		\item Misalkan $\phi: E \to L$ adalah $F$-pembenaman, $u \in E$ adalah unsur aljabar atas $F$, maka $\phi(u)$ juga merupakan unsur aljabar atas $F$ dan memiliki polinomial minimal yang sama dengan $u$.
		\item Jika unsur aljabar $u \in E$ dan $v \in L$ mempunyai polinomial minimal yang sama, yaitu $P$, maka terdapat tepat satu $F$-pembenaman $\iota: F(u) \to L$ sedemikian sehingga $\iota(u)=v$; pembenaman ini memenuhi $\Image(\iota)=F(v)$.
	\end{compactenum}
\end{proposition}
\begin{proof}
	Mari kita lihat bagian pertama. Misalkan $P \in F[X]$ adalah polinomial minimal dari $u$, karena $\phi$ adalah homomorfisme dari $F$-aljabar, maka $P(\phi(u)) = \phi(P(u)) = \phi(0) =0$, dengan demikian $\phi(u)$ adalah unsur aljabar; selanjutnya karena $P$ sudah tak tereduksi, itu pasti merupakan polinomial minimal dari $\phi(u)$.

	Sekarang perhatikan bagian kedua. Karena $u$ membangkitkan $F(u)$, $F$-pembenaman $\iota: F(u) \to L$ yang dicari sepenuhnya ditentukan oleh citra $v$ dari $u$, dan $\iota(F(u))=F(v)$. Tinggal membuktikan keberadaan $\iota$. Perhatikan diagram komutatif berikut:
	\begin{equation*}\begin{tikzcd}[row sep=tiny]
		F(u) \arrow[bend left=20, rr, "\iota"] & F[X]/(P) \arrow[l, "\sim"'] \arrow[r, "\sim"] & F(v) \arrow[hookrightarrow, r] & L \\
		u \arrow[phantom, u, sloped, "\in" description] & X+(P) \arrow[mapsto, l] \arrow[mapsto, r] \arrow[phantom, u, sloped, "\in" description] & v \arrow[phantom, u, sloped, "\in" description] &
	\end{tikzcd}\end{equation*}
\end{proof}

Dengan memanfaatkan Proposisi \ref{prop:adjunction-root}, untuk $F$ dan $P \in F[X]$ kita dapat secara bertahap membangun ekstensi berhingga $F = F_0 \hookrightarrow F_1 \hookrightarrow \cdots \hookrightarrow F_m = E$, dengan menambahkan satu akar $P \in F_i[X]$ setiap kali, sehingga $P$ akhirnya terurai menjadi hasil kali faktor linear di $E$. Dari perspektif teori, cara yang mengatasi semuanya sekaligus adalah membangun ekstensi yang cukup besar $\overline{F}|F$ sehingga setiap $P \in F[X]$ terurai menjadi hasil kali faktor linear di atas $\overline{F}$. Ini mengarah pada konsep medan tertutup secara aljabar dan penutup aljabar.

\begin{definition}\index{yu! tertutup secara aljabar (algebraically closed)}
	Jika setiap polinomial tak konstan $P \in E[X]$ di atas medan $E$ memiliki akar, atau secara ekuivalen memiliki faktor linear, maka $E$ disebut \emph{medan tertutup secara aljabar}.
\end{definition}
Contoh medan tertutup secara aljabar yang terkenal adalah medan bilangan kompleks $\CC$; inilah isi teorema dasar aljabar.

\begin{lemma}\label{prop:alg-closed-no-ext}
	Medan $E$ tertutup secara aljabar jika dan hanya jika ia tidak memiliki ekstensi aljabar nontrivial, yakni: $L|E$ merupakan ekstensi aljabar $\iff L=E$.
\end{lemma}
\begin{proof}
	Misalkan $E$ tertutup secara aljabar dan $L|E$ merupakan ekstensi aljabar. Untuk setiap $x \in L$, polinomial minimalnya $P_x$ atas $E$ pasti mempunyai faktor linear; karena $P_x$ tak tereduksi, maka $\deg P_x = 1$, sehingga $x \in E$.
	
	Sebaliknya, andaikan $E$ tidak tertutup secara aljabar dan $P \in E[X]$ tidak memiliki akar di $E$; kita boleh menganggap $P$ tak tereduksi. Proposisi \ref{prop:adjunction-root} kemudian memberikan ekstensi medan $L := E[X]/(P)$ dari $E$ dengan $[L:E] = \deg P > 1$. Terbukti.
\end{proof}

\begin{lemma}\label{prop:alg-closed-decomp}
	Misalkan $E|F$ suatu ekstensi aljabar. Maka $E$ tertutup secara aljabar jika dan hanya jika setiap polinomial tak konstan $P \in F[X]$ terurai menjadi hasil kali faktor-faktor linear di $E$.
\end{lemma}
\begin{proof}
	Keniscayaan langsung mengikuti definisi medan tertutup secara aljabar; sekarang buktikan kecukupannya. Berdasarkan Lemma \ref{prop:alg-closed-no-ext}, cukup dibuktikan $L=E$ untuk setiap ekstensi aljabar $L|E$. Proposisi \ref{prop:alg-ext-dist} memastikan bahwa $L|F$ juga aljabar. Untuk sembarang $x \in L$, menurut asumsi polinomial minimal $P_x \in F[X]$ terurai di $E[X]$ sebagai $P_x = \prod_{i=1}^n (X - \alpha_i)$. Karena $P_x(x)=0$, terdapat $1 \leq i \leq n$ dengan $x = \alpha_i \in E$. Terbukti.
\end{proof}
Syarat dari lemmata di atas juga dapat diperbaiki menjadi: Setiap polinomial bukan konstanta $P \in F[X]$ memiliki akar di $E$. Kita akan membuktikannya dalam proposisi \ref{prop:alg-closed-root} menggunakan teori Galois.

\begin{definition}\index{daishubibao@penutup aljabar (algebraic closure)}
	Ekstensi medan $\overline{F}|F$ yang memenuhi
	\begin{inparaenum}[(i)]
		\item $\overline{F}$ adalah medan tertutup secara aljabar,
		\item $\overline{F}|F$ adalah ekstensi aljabar,
	\end{inparaenum}
	disebut \emph{penutup aljabar} dari $F$.
\end{definition}
Proposisi \ref{prop:alg-ext-dist} segera mengarah pada sebuah sifat umum: Misalkan $E|F$ adalah ekstensi aljabar, dan $\overline{E}|E$ adalah penutup aljabar dari $E$, maka $\overline{E}|F$ juga merupakan penutup aljabar dari $F$.

\begin{theorem}[E.\ Steinitz]\label{prop:alg-closure}
	Untuk setiap medan $F$, penutup aljabar $\overline{F}|F$ ada, dan unik dalam arti $F$-isomorfik.
\end{theorem}
\begin{proof}[E.\ Artin]
	Pertama buktikan keunikan. Misalkan $\overline{F}|F$ dan $\overline{F}'|F$ adalah penutup aljabar. Definisikan $\mathcal{P}$ sebagai himpunan semua data $(E, \iota)$, dengan $E|F$ suatu subekstensi dari $\overline{F}|F$ dan $\iota \in \Hom_F(E, \overline{F}')$. Jelas $\left( F, F \hookrightarrow \overline{F}' \right) \in \mathcal{P}$, sehingga $\mathcal{P}$ tak kosong. Beri $\mathcal{P}$ urutan parsial
	\[ (E, \iota) \leq (E_1, \iota_1) \iff (E \subset E_1) \wedge (\iota_1|_E = \iota). \]
	Mudah dilihat bahwa dalam $(\mathcal{P}, \leq)$, setiap rantai $\{ (L_i, \iota_i)\}_{i \in I}$ mempunyai batas atas $(L, \iota)$: ambil subekstensi $L := \bigcup_i L_i$ dan $\iota: L \to \overline{F}'$ sedemikian sehingga $\iota|_{L_i} = \iota_i$. Lemma Zorn (Teorema \ref{prop:Zorn}) menjamin bahwa $(\mathcal{P}, \leq)$ mempunyai unsur maksimal $(E, \phi)$.
	
	Kita nyatakan $E=\overline{F}$. Jika tidak, ambil $u \in \overline{F} \smallsetminus E$, dengan polinomial minimal $P_u \in E[X]$ atas $E$. Melalui $\phi$, pandang $\overline{F}'$ sebagai ekstensi medan dari $E$. Maka $P_u$ mempunyai akar $v$ di $\overline{F}'$, sehingga Proposisi \ref{prop:field-embedding} memberikan pembenaman-$E$ yang unik $E(u) \hookrightarrow \overline{F}'$ dan memetakan $u$ ke $v$; ini bertentangan dengan maksimalitas $\phi$.
	
	Tinggal membuktikan $\phi(\overline{F})=\overline{F}'$ untuk memperoleh $F$-isomorfisme $\phi: \overline{F} \rightiso \overline{F}'$. Karena $\overline{F}$ tertutup secara aljabar, citra isomorfiknya $\phi(\overline{F})$ juga demikian. Karena $\overline{F}'|F$ aljabar, $\overline{F}'|\phi(\overline{F})$ juga aljabar; Lemma \ref{prop:alg-closed-no-ext} lalu memberikan $\phi(\overline{F})=\overline{F}'$. Keunikan terbukti.
	
	Kita berikan dua argumen untuk bagian eksistensi. Argumen pertama ringkas, tetapi memerlukan konstruksi produk tensor $F$-aljabar yang diperkenalkan di \S\ref{sec:algebra-tensor-product}. Untuk setiap polinomial tak tereduksi $P \in F[X]$, pembahasan sebelumnya dan Proposisi \ref{prop:adjunction-root} memungkinkan kita membangun secara bertahap ekstensi berhingga $E_P|F$ sehingga $P$ terurai menjadi faktor-faktor linear di atas $E_P$. Bagaimana menemukan ekstensi medan yang dibangkitkan oleh semua $E_P$, atau setidaknya suatu $F$-aljabar komutatif semacam itu? Sudut pandang kategoris memberikan petunjuk: perhatikan produk tensor tak hingga dari $F$-aljabar komutatif (Definisi--Teorema \ref{def:inf-tensor-product})
	\[ A := \bigotimes_{\substack{P \in F[X] \\ \text{tak tereduksi monik } }} E_P. \]
	Kita mempunyai keluarga $F$-homomorfisme $E_P \hookrightarrow A$ yang citranya membangkitkan $F$-aljabar komutatif $A$; Catatan \ref{rem:algebra-otimes-zero} menjamin $A$ tak nol. Proposisi \ref{prop:existence-maximal-ideal}, yang bertumpu pada aksioma pilihan, menyatakan bahwa $A$ mempunyai ideal maksimal $I$. Karena itu aljabar hasil bagi $\overline{F} := A/I$ adalah medan dan tetap dibangkitkan oleh citra suatu keluarga $F$-pembenaman $E_P \to \overline{F}$. Jadi $\overline{F}|F$ adalah ekstensi aljabar tempat setiap $P$ terurai menjadi faktor linear. Proposisi \ref{prop:alg-closed-decomp} langsung menyiratkan bahwa $\overline{F}$ tertutup secara aljabar.
	
	Argumen kedua adalah secara langsung menggunakan teorema \ref{prop:Zorn} untuk mencari sebuah ekstensi aljabar maksimum dari $F$. Ambil himpunan $\Omega$ yang cukup besar sehingga $|\Omega| > \max\{|F|, \aleph_0\}$, dan $F \subset \Omega$; tujuan kita adalah menyisipkan semua ekstensi aljabar sebagai subset dari himpunan 'alam semesta' $\Omega$. Definisikan himpunan $\mathcal{Q}$ sebagai semua ekstensi aljabar $E|F$, yang sesuai untuk
	\begin{inparaenum}[(i)]
		\item sebagai himpunan memiliki $F \subset E \subset \Omega$,
		\item ekstensi medan $E|F$ ditentukan oleh peta inklusi $F \hookrightarrow E$.
	\end{inparaenum}
	Mudah terlihat bahwa $\mathcal{Q}$ memang merupakan himpunan, $F|F \in \mathcal{Q}$. Selanjutnya memberikan $\mathcal{Q}$ urutan parsial
	\[ E|F \leq E_1|F \iff E \subset E_1 \quad (\text{submedan}). \]
	Konstruksi gabungan yang biasa menunjukkan bahwa setiap rantai dalam $(\mathcal{Q}, \leq)$ mempunyai batas atas, sehingga terdapat unsur maksimal $\overline{F}|F$. Kita nyatakan bahwa setiap ekstensi aljabar dari $\overline{F}$, yaitu $L|\overline{F}$, harus memenuhi $L=\overline{F}$; Lemma \ref{prop:alg-closed-no-ext} lalu menyiratkan bahwa $\overline{F}$ tertutup secara aljabar. Mula-mula kita harus menyesuaikan $L|F$ agar menjadi unsur $\mathcal{Q}$.
	
	Proposisi \ref{prop:alg-ext-cardinality} menyatakan bahwa untuk setiap ekstensi aljabar $E|F$ terdapat $|E| \leq \max\{|F|, \aleph_0\} < |\Omega|$, tentu saja ini berlaku untuk $E = L, \overline{F}$. Kita memerlukan sedikit estimasi kardinalitas:
	\begin{itemize}
		\item Menurut pemilihan $\Omega$, $|\Omega \smallsetminus \overline{F}|$ dan $|\overline{F}|$ keduanya tidak nol, dan salah satunya tak hingga; sebenarnya $|\Omega \smallsetminus \overline{F}| \geq |\overline{F}|$, jika tidak maka korolari \ref{prop:cardinal-max} akan menghasilkan kontradiksi $|\Omega| = \max\{|\Omega \smallsetminus \overline{F}|, |\overline{F}|\} = |\overline{F}| < |\Omega|$.
		\item Maka $|\Omega| = \max\{|\Omega \smallsetminus \overline{F}|, |\overline{F}|\} = |\Omega \smallsetminus \overline{F}|$, selanjutnya
			\[ |L \smallsetminus \overline{F}| \leq |L| \leq \max\{|F|, \aleph_0\} < |\Omega| = |\Omega \smallsetminus \overline{F}|. \] 
		\item Oleh karena itu ada injeksi dari himpunan $f: L \hookrightarrow \Omega$ sehingga $f|_{\overline{F}} = \identity_{\overline{F}}$; memindahkan struktur medan menggunakan $f$ ke $f(L)$, tanpa mengubah struktur medan dari $\overline{F} \subset \Omega$, ini membuat $f(L)|F$ menjadi anggota dari $\mathcal{Q}$.
	\end{itemize}
	Dengan demikian, dalam $(\mathcal{Q}, \leq)$ berlaku $f(L)|F \geq \overline{F}|F$. Maksimalitas memberikan $f(L)=\overline{F}$, sehingga $L=\overline{F}$. Terbukti.
\end{proof}

\begin{corollary}\label{prop:alg-ext-embedding}
	Misalkan $\overline{F}|F$ penutup aljabar dan $E|F$ ekstensi aljabar. Maka terdapat $F$-pembenaman $\iota \in \Hom_F(E, \overline{F})$. Jika $E|F$ berhingga, maka $|\Hom_F(E, \overline{F})| \leq [E:F]$.
\end{corollary}
\begin{proof}
	Ambil penutup aljabar $\overline{E}$ dari $E$ dan pandang $E$ sebagai submedannya. Menurut definisi, $\overline{E}|F$ juga merupakan penutup aljabar; keunikan memberikan $F$-isomorfisme $\phi: \overline{E} \rightiso \overline{F}$. Ambil $\iota := \phi|_E$ untuk membuktikan pernyataan pertama. Sekarang andaikan $E = F(x_1, \ldots, x_n)$ ekstensi berhingga dari $F$, dan perhatikan peta pembatasan
	\begin{multline*}
		\Hom_F(F(x_1, \ldots, x_n), \overline{F}) \xrightarrow{\text{res}_n} \Hom_F(F(x_1, \ldots, x_{n-1}), \overline{F}) \xrightarrow{\text{res}_{n-1}} \\
		\cdots \xrightarrow{\text{res}_1} \Hom_F(F, \overline{F}) = \{\ast\}.
	\end{multline*}
	Untuk sembarang $\phi \in \Hom_F(F(x_1, \ldots, x_{i-1}), \overline{F})$, Proposisi \ref{prop:field-embedding} memastikan bahwa ukuran serat $\text{res}_i^{-1}(\phi)$ tidak melebihi derajat polinomial minimal $x_i$ atas $F(x_1, \ldots, x_{i-1})$. Karena itu
	\[ \left| \Hom_F(F(x_1, \ldots, x_n), \overline{F}) \right| \leq \prod_{i=1}^n [F(x_1, \ldots, x_i): F(x_1, \ldots, x_{i-1})] = [E:F], \]
	Terbukti.
\end{proof}

\begin{proposition}\label{prop:alg-closure-sub}
	Jika $\Omega$ tertutup secara aljabar dan $\Omega|F$ ekstensi medan, tetapkan $\overline{F} := \{ x \in \Omega: \text{aljabar atas $F$} \}$. Maka $\overline{F}|F$ adalah penutup aljabar dari $F$.
\end{proposition}
\begin{proof}
	Setiap $x \in \overline{F}$ menghasilkan ekstensi berhingga $F(x)=F[x]$, dan karena ekstensi aljabar tertutup terhadap komposisi, maka $\overline{F} = \bigvee_{x \in \overline{F}} F(x)$ memang merupakan ekstensi aljabar dari $F$. Setiap polinomial tak tereduksi monik $P \in F[X]$ dapat dipecah di $\Omega$ menjadi faktor-faktor linear $\prod_\alpha (X-\alpha)$; menurut definisi, $\alpha \in \overline{F}$, sehingga Lemma \ref{prop:alg-closed-decomp} menyiratkan bahwa $\overline{F}$ adalah medan tertutup secara aljabar.
\end{proof}

Dari ketertutupan aljabar $\CC$, contoh pertama penutup aljabar adalah $\CC|\R$. Contoh berikutnya ialah, di dalam $\CC$, himpunan semua bilangan aljabar atas $\Q$, yang membentuk submedan $\overline{\Q}$; Proposisi \ref{prop:alg-closure-sub} menyatakan bahwa $\overline{\Q}$ adalah penutup aljabar dari $\Q$. Dalam kasus umum, Teorema \ref{prop:alg-closure} hanya memberikan keberadaan dan keunikan penutup aljabar $\overline{F}|F$ secara abstrak. Setiap jalur pembuktian tidak dapat menghindari aksioma pilihan dan karenanya bersifat nonkonstruktif.

\section{Medan Pemisah dan Ekstensi Normal}\label{sec:splitting-field}
Kita sudah melihat bahwa setiap $P \in F[X]$ yang bukan konstanta terurai menjadi faktor linear pada beberapa ekstensi berhingga $E|F$. Bagian ini bertujuan untuk memberikan penjelasan yang lebih rinci tentang konstruksi semacam ini.

\begin{definition}\index{fenlieyu@medan pemisah (splitting field)}
	Misalkan $\mathcal{P}$ adalah satu keluarga polinomial tidak konstan di $F[X]$. Jika ekstensi medan $E|F$ memenuhi
	\begin{compactitem}
		\item Setiap $P \in \mathcal{P}$ terurai di $E$ menjadi faktor-faktor linear, yakni $P = c_P \prod_{j=1}^{n_P}(X - \alpha_{P,j})$, dengan $\alpha_{P,j} \in E$ dan $c_P \in F^\times$;
		\item akar-akar $\left\{ \alpha_{P,j} : P \in \mathcal{P},\; 1 \leq j \leq n_P \right\}$ membangkitkan $E$ di atas $F$.
	\end{compactitem}
	maka $E|F$ bagi keluarga polinomial $\mathcal{P}$ disebut \emph{medan pemisah}.
\end{definition}
Medan pemisah selalu merupakan ekstensi aljabar. Kasus yang dapat dibayangkan paling rumit adalah $\mathcal{P} := F[X] \smallsetminus F$; dalam hal ini Lemma \ref{prop:alg-closed-decomp} menunjukkan bahwa medan pemisah tepat merupakan penutup aljabar dari $F$, yang keberadaannya sudah dibuktikan. Namun, dalam kajian persamaan polinomial kita lebih sering mempertimbangkan medan pemisah satu polinomial $P \in F[X]$.

\begin{proposition}
	Untuk setiap keluarga polinomial tak konstan $\mathcal{P}$, medan pemisah $E|F$ selalu ada dan unik hingga $F$-isomorfisme. Medan pemisah suatu polinomial $P$ berderajat $n \geq 1$ atas $F$ mempunyai derajat $ \leq n!$.
\end{proposition}
\begin{proof}
	Pertama buktikan keberadaan. Ambil penutup aljabar $\overline{F}|F$. Setiap $P \in \mathcal{P}$ terurai di $\overline{F}$ sebagai $P = c_P \prod_{j=1}^{n_P}(X - \alpha_{P,j})$. Subekstensi yang dibangkitkan oleh semua akar, yaitu $E = F\left( \alpha_{P,j} : P \in \mathcal{P},\; 1 \leq j \leq n_P \right)$, adalah yang dicari. Di antara subekstensi $\overline{F}|F$, jelas inilah satu-satunya pilihan untuk medan pemisah $\mathcal{P}$.

	Sekarang buktikan keunikan. Misalkan $E|F$ dan $E'|F$ medan pemisah $\mathcal{P}$, dan ambil masing-masing penutup aljabar $\overline{F}|E$ dan $\overline{F}'|E'$. Kita boleh memandang $E \subset \overline{F}$ dan $E' \subset \overline{F}'$. Menurut definisi, $\overline{F}|F$ dan $\overline{F}'|F$ juga penutup aljabar; Teorema \ref{prop:alg-closure} memberikan $F$-isomorfisme $u: \overline{F} \rightiso \overline{F}'$. Karena $E|F$ dibangkitkan oleh semua akar polinomial dalam $\mathcal{P}$ yang berada di $\overline{F}$, sifat yang sama dibawa oleh $u$ ke $u(E) \subset \overline{F}'$. Jadi $u(E)|F$ sama dengan $E'|F$, dan $F$-isomorfisme yang dicari adalah $u|_E$.

	Untuk derajatnya, andaikan $\deg P = n \geq 1$, dan tulis akar-akar $P$ dalam medan pemisah $E$ sebagai $\alpha_1, \ldots, \alpha_n$. Jika $n=1$, jelas $E=F$. Jika $n>1$, tetapkan $F_1 := F(\alpha_1)$. Karena $[F_1:F] \leq \deg P = n$ dan $E$ adalah medan pemisah $P/(X-\alpha_1)$ atas $F_1$, Proposisi \ref{prop:field-tower-degree} secara rekursif memberikan $[E:F] = [E:F_1][F_1:F] \leq (n-1)! n = n!$.
\end{proof}

\begin{lemma}\label{prop:alg-ext-endomorphism}
	Untuk setiap ekstensi aljabar $L|F$, setiap $F$-pembenaman $\iota: L \to L$ adalah isomorfisme. Dengan kata lain, $\End_F(L)=\Aut_F(L)$.
\end{lemma}
\begin{proof}
	Cukup dibuktikan bahwa $\iota$ surjektif. Untuk sembarang $y \in L$, tulis akar-akar polinomial minimal $P_y$ yang berada di $L$ sebagai $y_1, \ldots, y_m$, dan definisikan subekstensi $L_0 := F(y_1, \ldots, y_m)$ dari $L|F$. Karena $F$-pembenaman $\iota$ menginduksi permutasi pada himpunan akar $\{y_1, \ldots, y_m\}$, berlaku $\iota|_{L_0}: L_0 \hookrightarrow L_0$. Lemma \ref{prop:finite-vs-algebraic} menyatakan bahwa $L_0|F$ berhingga; berdasarkan dimensi, $\iota|_{L_0}$ merupakan isomorfisme. Khususnya, $y \in \Image(\iota)$.
\end{proof}

\begin{definition-theorem}\label{def:normal-ext}\index{yukuozhang!normal (normal)}
	Untuk ekstensi aljabar $E|F$, sifat-sifat berikut setara:
	\begin{enumerate}[\bfseries {N}.1]
		\item Jika polinomial tak tereduksi $P \in F[X]$ mempunyai akar di $E$, maka polinomial itu terurai di $E[X]$ sebagai hasil kali faktor-faktor linear.
		\item Tetapkan penutup aljabar $\overline{F}|E$ dan pandang $E$ sebagai submedan $\overline{F}$. Maka setiap $\iota \in \Hom_F(E, \overline{F})$ memenuhi $\iota(E)=E$.
		\item Ada suatu keluarga polinomial bukan konstan $\mathcal{P}$ sehingga $E|F$ adalah medan pemisah $\mathcal{P}$.
	\end{enumerate}
	Ekstensi aljabar yang memenuhi salah satu kondisi di atas disebut \emph{ekstensi normal}.
\end{definition-theorem}
\begin{proof}
	(\textbf{N.1}) $\implies$ (\textbf{N.2}): Diberikan $\iota \in \Hom_F(E, \overline{F})$. Tuliskan polinomial minimal dari $x \in E$ sebagai $P_x \in F[X]$; maka $\iota(x)$ tetap merupakan akar dari $P_x$. Menurut syarat, $P_x$ terurai dalam $E[X]$ menjadi faktor-faktor linear $(X-\alpha_1) \cdots (X-\alpha_n)$, sehingga $\exists i,\; \iota(x) = \alpha_i \in E$. Karena $x$ dapat dipilih sembarang, langsung diperoleh $\iota(E) \subset E$. Dengan menerapkan Lemma \ref{prop:alg-ext-endomorphism}, diperoleh $\iota(E)=E$.
	
	(\textbf{N.2}) $\implies$ (\textbf{N.1}): Jika $P$ mempunyai akar $x \in E$, ambil sembarang akar $y \in \overline{F}$ dari $P$. Proposisi \ref{prop:field-embedding} memberikan $F$-pembenaman $\iota_0: F(x) \to \overline{F}$ dengan $\iota_0(x)=y$. Melalui $\iota_0$, $\overline{F}|F(x)$ dapat dipandang sebagai penutup aljabar. Korolari \ref{prop:alg-ext-embedding} memberikan diagram komutatif berikut
	\[\begin{tikzcd}[row sep=small, column sep=small]
		E \arrow[rr, "\exists\; \iota"] & & \overline{F} \\
		& F(x) \arrow[hookrightarrow, lu, "\text{inklusi}"] \arrow[ru, "\iota_0"'] &
	\end{tikzcd} \quad \text{maka}\; \iota \in \Hom_F(E, \overline{F}). \]
	Oleh karena itu $y = \iota_0(x) = \iota(x) \in E$, terlihat bahwa semua akar dari $P$ berada di $E$.
	
	(\textbf{N.1}) $\implies$ (\textbf{N.3}): Ambil $\mathcal{P}$ sebagai semua polinomial tak tereduksi yang memiliki akar di $E$, kondisi tersebut mengimplikasikan medan pemisah $\subset E$. Karena setiap $x \in E$ adalah akar dari polinomial minimal $P_x$, medan pemisah yang sesuai adalah $E$.
	
	(\textbf{N.3}) $\implies$ (\textbf{N.2}): Untuk setiap $\iota \in \Hom_F(E, \overline{F})$ dan $P \in \mathcal{P}$, pembenaman $\iota$ menginduksi permutasi pada $R_P := \{x \in \overline{F}: P(x)=0 \}$. Medan pemisah $E$ dibangkitkan di atas $F$ oleh $\bigcup_{P \in \mathcal{P}} R_P$, sehingga $\iota(E)=E$.
\end{proof}

Hal ini secara alami mengantar pada istilah berikut.
\begin{definition}\index{gonge}
	Dua subekstensi $E|F$, $E'|F$ dari $\Omega|F$ disebut konjugat jika terdapat $\sigma \in \Aut_F(\Omega)$ dengan $\sigma(E)=E'$. Jika untuk $x, y \in \Omega$ terdapat $\sigma \in \Aut_F(\Omega)$ dengan $\sigma(x)=y$, maka $x$ dan $y$ disebut konjugat.
\end{definition}
Konjugasi unsur jelas merupakan relasi ekuivalensi. Salah satu asal istilah ini ialah kasus $\Omega|F = \CC|\R$: karena $\CC = \R \oplus \R i \simeq \R[X]/(X^2+1)$ dan akar-akar $X^2+1$ adalah $\pm i$, automorfisme $\sigma \in \Aut_{\R}(\CC)$ ditentukan secara unik oleh citra $\sigma(i) = \pm i$. Jadi $\sigma$ harus berupa $\identity_{\CC}$ atau konjugasi kompleks $x + yi \mapsto x - yi$.

\begin{proposition}\label{prop:normal-ext-prolongation}
	Misalkan $L|F$ ekstensi normal dan $E|F$ subekstensinya. Maka setiap $\iota \in \Hom_F(E, L)$ dapat diperluas menjadi suatu $\tilde{\iota} \in \Aut_F(L)$.
\end{proposition}
\begin{proof}
	Menurut Lemma \ref{prop:alg-ext-endomorphism}, cukup memperpanjang $\iota$ menjadi $\tilde{\iota}: L \to L$. Ambil penutup aljabar $\overline{F}$ dari $L$, dan tetap tulis komposit $E \xrightarrow{\iota} L \hookrightarrow \overline{F}$ sebagai $\iota$. Maka $\iota: E \to \overline{F}$ memberi penutup aljabar dari $E$ dan juga dari $F$, sehingga notasi $\overline{F}$ wajar. Pada penutup aljabar $E \xrightarrow{\iota} \overline{F}$ dan inklusi $E \hookrightarrow L$, terapkan Korolari \ref{prop:alg-ext-embedding}; diperoleh pembenaman medan $\tilde{\iota}: L \to \overline{F}$ dengan $\tilde{\iota}|_E = \iota$. Tinggal dibuktikan $\tilde{\iota}(L) \subset L$.

	Karena $\overline{F}$ melalui $\iota$ menjadi penutup aljabar dari $F$, kenormalan $L|F$ dan \textbf{N.2} langsung memberikan $\overline{\iota}(L) \subset L$, sebagaimana hendak dibuktikan.
%	Pertama, pertimbangkan kasus $L=\overline{F}$ merupakan penutup aljabar dari $F$. Maka $\iota: E \hookrightarrow \overline{F}$ juga memberikan penutup aljabar dari $E$. Korolari \ref{prop:alg-ext-embedding} menyatakan bahwa diagram berikut dapat dilengkapi dengan panah putus-putus $\tilde{\iota}$ sehingga komutatif
%	\[\begin{tikzcd}[column sep=small, row sep=small]
%		\overline{F} \arrow[dashed, rr, "\exists\; \tilde{\iota}"] & & \overline{F} \\
%		& E \arrow[hookrightarrow, lu] \arrow[ru, "\iota"'] &
%	\end{tikzcd}\]
%	Inilah $\tilde{\iota}$ yang diperlukan. Untuk $L$ umum, ambil penutup aljabarnya $\overline{F}$, yang juga merupakan penutup aljabar dari $F$; menurut langkah sebelumnya, perpanjang komposisi $E \xrightarrow{\iota} L \hookrightarrow \overline{F}$ menjadi $\overline{\iota}: \overline{F} \to \overline{F}$. Dari kenormalan $L$ dan \textbf{N.2} diperoleh $\overline{\iota}|_L: L \to L$.
\end{proof}

Ini memberikan generalisasi yang berguna dari kondisi \textbf{N.2} dalam definisi ekstensi normal.
\begin{corollary}\label{prop:normal-subext}
	Untuk ekstensi normal $L|F$, subekstensi $E|F$ normal jika dan hanya jika:
	\begin{compactdesc}
		\item[\bfseries {N}.2$'$] Untuk setiap $\sigma \in \Aut_F(L)$ berlaku $\sigma(E) \subset E$.
	\end{compactdesc}
	Sebagai aplikasi, jika dalam menara empat tingkat dari ekstensi aljabar $L|E|K|F$, $L|F$ dan $E|F$ keduanya normal, maka $\Hom_F(K, L) = \Hom_F(K, E)$.
\end{corollary}
\begin{proof}
	Ambil penutup aljabar $\overline{F}|L$. Jika $E|F$ normal, maka karena $\Aut_F(L) \subset \Hom_F(L, \overline{F})$, menurut \textbf{N.2} pasti berlaku $\sigma(E) \subset E$. Sebaliknya, misalkan $\forall\sigma\; \sigma(E) \subset E$, proposisi \ref{prop:normal-ext-prolongation} menjamin bahwa setiap $\iota \in \Hom_F(E, \overline{F})$ selalu dapat diperpanjang ke $\overline{F} \to \overline{F}$, yang pembatasannya pada $L$ dicatat sebagai $\sigma$; kondisi \textbf{N.2} menjamin $\sigma \in \End_F(L) = \Aut_F(L)$, sehingga $\iota(E) = \sigma(E) \subset E$. Ini membuktikan bagian pertama.

	Mengenai bagian kedua, jelas $\Hom_F(K,E) \subset \Hom_F(K,L)$. Jika $\iota \in \Hom_F(K,L)$, gunakan proposisi \ref{prop:normal-ext-prolongation} untuk memperluas $\iota$ menjadi $\sigma \in \Aut_F(L)$, kemudian dari bagian pertama diturunkan $\iota(K) \subset \sigma(E) \subset E$. Terbukti.
\end{proof}
Dalam \textbf{N.2$'$}, ambil $L = \overline{F}$; langsung terlihat bahwa ekstensi normal tidak lain adalah subekstensi $\overline{F}|F$ yang tertutup terhadap relasi konjugasi.

\begin{corollary}\label{prop:root-conjugate}
	Di dalam penutup aljabar $\overline{F}|F$, dua unsur $x$, $y$ konjugat jika dan hanya jika mempunyai polinomial minimal yang sama.
\end{corollary}
\begin{proof}
	Jelas bahwa $F$-automorfisme mempertahankan polinomial minimal unsur. Sebaliknya, andaikan $x,y$ sama-sama akar polinomial tak tereduksi $P \in F[X]$. Proposisi \ref{prop:field-embedding} memberikan $F$-pembenaman $\iota: F(x) \to \overline{F}$ dengan $\iota(x)=y$. Proposisi \ref{prop:normal-ext-prolongation} memperluas $\iota$ menjadi $\sigma \in \Aut_F(\overline{F})$, sehingga $\sigma(x)=y$ dan keduanya konjugat.
\end{proof}

Meskipun subekstensi suatu ekstensi normal belum tentu normal, ekstensi normal tetap mempunyai sifat ketertutupan yang baik terhadap operasi lain.

\begin{proposition}\label{prop:normal-ext-properties}
	Ekstensi normal memenuhi sifat-sifat berikut.
	\begin{enumerate}[(i)]
		\item Ekstensi $L|F$ normal $\implies$ untuk setiap medan antara $E$, ekstensi $L|E$ normal.
		\item Misalkan $L|F$, $M|F$ adalah subekstensi dari $\Omega|F$, maka $L|F$ normal $\implies$ $LM|M$ normal.
		\item Dalam ekstensi tetap $\Omega|F$, kompositum setiap keluarga subekstensi normal dan irisan tak kosong setiap keluarga tersebut tetap normal.
	\end{enumerate}
\end{proposition}
\begin{proof}
	Misalkan $L|F$ dan $M|F$ keduanya adalah subekstensi dari $\Omega|F$, jika $L|F$ adalah medan pemisah dari $\mathcal{P} \subset F[X]$, maka $LM|M$ adalah medan pemisah dari $\mathcal{P} \subset M[X]$, sehingga merupakan ekstensi normal. Ini membuktikan (ii); dengan mengambil kasus khusus $\Omega=L$ maka $M=E \subset L$ didapatkan (i).
	
	Misalkan $\{L_i|F\}_{i \in I}$ keluarga subekstensi normal dari $\Omega|F$, dengan $I$ tak kosong. Irisan $\bigcap_i L_i|F$ dan kompositum $\bigvee_i L_i | F$ tetap merupakan ekstensi aljabar (Proposisi \ref{prop:alg-ext-dist}). Untuk $\bigvee_i L_i$, ambil penutup aljabar $\overline{F}$ yang memuatnya; ia sekaligus merupakan penutup aljabar dari $\bigcap_i L_i$, setiap $L_i$, dan $F$, sedangkan $\bigvee_i L_i$ dapat dipandang sebagai kompositum di dalam $\overline{F}|F$. Pertama buktikan bahwa $\bigvee_i L_i|F$ normal: untuk setiap $\iota \in \Hom_F\left(\bigvee_i L_i , \overline{F}\right)$, normalitas memberikan
	\begin{gather*}
		\iota\left( \bigvee_{i \in I} L_i \right) = \bigvee_{i \in I} \iota(L_i) = \bigvee_{i \in I} L_i.
	\end{gather*}
	Selanjutnya buktikan bahwa $\bigcap_i L_i|F$ normal. Karena $\bigvee_i L_i \big| \bigcap_i L_i$ aljabar, Korolari \ref{prop:alg-ext-embedding} menyatakan bahwa setiap $\iota \in \Hom_F\left(\bigcap_i L_i , \overline{F}\right)$ dapat diperluas menjadi $\bigvee_i L_i \to \overline{F}$, lalu dibatasi pada setiap $L_i$; normalitas memberikan
	\begin{gather*}
		\iota\left( \bigcap_{i \in I} L_i \right) = \bigcap_{i \in I} \iota(L_i) = \bigcap_{i \in I} L_i.
	\end{gather*}
	Terbukti.
\end{proof}

\begin{definition}\label{def:normal-closure}\index{zhengguibibao@penutup normal (normal closure)}
	Misalkan $E|F$ ekstensi aljabar dan $\overline{F}|E$ penutup aljabar yang dipilih. \emph{Penutup normal} $M|F$ dari $E|F$ didefinisikan sebagai irisan semua subekstensi normal dalam $\overline{F}|F$ yang memuat $E|F$.
\end{definition}
Karena $\overline{F}|F$ normal, irisan ini tak kosong. Proposisi \ref{prop:normal-ext-properties} menunjukkan bahwa $M|F$ adalah subekstensi normal terkecil dalam $\overline{F}|F$ yang memuat $E$. Selanjutnya,
\[ M \xlongequal{\because\; \textbf{N.2$'$}} \bigvee_{\sigma \in \Aut_F(\overline{F})} \sigma(E) \xlongequal{\because \text{proposisi \ref{prop:normal-ext-prolongation}}} \bigvee_{\sigma \in \Hom_F(E, \overline{F})} \sigma(E). \]
Jika lebih lanjut $E|F$ berhingga, Korolari \ref{prop:alg-ext-embedding} menyiratkan bahwa kompositum di atas hanya mempunyai berhingga banyak suku; jadi $M|F$ juga ekstensi berhingga. Perhatikan contoh standar berikut.

\begin{example}\label{eg:x3-2-splitting}
	Di dalam $\overline{\Q}|\Q$, subekstensi $\Q(2^{1/3})|\Q$ tidak normal, sebab $\Q(2^{1/3}) \subset \R$, sedangkan di dalam $\overline{\Q}$ masih terdapat konjugat $\omega 2^{1/3}$ dan $\omega^2 2^{1/3}$ dari $2^{1/3}$, dengan $\omega = e^{2\pi i/3}$; ketiganya adalah akar polinomial minimal $X^3 - 2$ dalam $\overline{\Q}$. Penutup normalnya ialah $\Q(2^{1/3},\omega)$.
\end{example}

\section{Separabilitas}
Untuk setiap medan $F$, Proposisi \ref{prop:polynomial-derivation}, melalui rumus biasa $(\sum_{k \geq 0} a_k X^k)' = \sum_{k \geq 1} ka_k X^{k-1}$, mendefinisikan operasi turunan $P \mapsto P'$ pada $F[X]$. Berikutnya tetapkan penutup aljabar $\overline{F}|F$. Untuk polinomial tak nol $P, Q$, tulis faktor persekutuan monik terbesarnya sebagai $(P, Q)$; faktor ini ditentukan secara unik.

\begin{lemma}\label{prop:coprime-basechange}
	Misalkan $P,Q \in F[X]$ tidak nol, $L|F$ adalah ekstensi medan sembarang, maka $(P,Q)=1$ berlaku dalam $F[X]$ jika dan hanya jika berlaku dalam $L[X]$.
\end{lemma}
\begin{proof}
	Jika $P,Q$ memiliki faktor umum yang bukan konstanta di $F[X]$, maka hal yang sama berlaku di $L[X]$. Jika keduanya relatif prima di $F[X]$, maka ada $U,V \in F[X]$ sedemikian sehingga $UP+VQ=1$, dan persamaan ini di $L[X]$ menyiratkan bahwa $P,Q$ relatif prima.
\end{proof}

\begin{lemma}\label{prop:separable-polynomial-gen}
	Syarat perlu dan cukup agar polinomial tidak nol $P \in F[X]$ memiliki akar ganda di medan pemisahnya adalah $(P,P') \neq 1$.
\end{lemma}
\begin{proof}
	Lemma \ref{prop:coprime-basechange} menunjukkan bahwa kondisi $(P, P') \neq 1$ adalah ekuivalen di $F$ dan di medan pemisah $L$. Semua dapat direduksi menjadi persamaan yang sudah dikenal berikut: untuk setiap ekstensi medan $L|F$ dan $c \in L$, terdapat $Q \in L[X]$ sehingga
	\begin{gather}\label{eqn:derivation-linear-approx}
		P = P(c) + (X-c) P'(c) + (X-c)^2 Q.
	\end{gather}
	Memang, karena $X \mapsto X-c$ menginduksi $L[X]$ sebagai automorfisme $L$-aljabar, masalah ini mudah direduksi ke kasus yang jelas dari $c=0$.
\end{proof}
Sebagai tambahan, cara lain untuk memeriksa akar ganda ialah menggunakan diskriminan yang diperkenalkan dalam Contoh \ref{eg:polynomial-discriminant}.

\begin{theorem}\label{prop:separable-polynomial}
	Untuk polinomial tak tereduksi $P \in F[X]$, pernyataan berikut setara.
	\begin{enumerate}[(i)]
		\item $P$ memiliki akar ganda di penutup aljabar $\overline{F}$,
		\item $P$ memiliki akar ganda pada medan pemisah,
		\item $P'=0$,
		\item $\mathrm{char}(F)=p > 0$, dan $P$ berbentuk seperti $\sum_{k=0}^n a_k X^{pk}$.
	\end{enumerate}
\end{theorem}
\begin{proof}
	(i) $\implies$ (ii) jelas.

	(ii) $\implies$ (iii): Jika $P$ mempunyai akar ganda dalam medan pemisahnya, maka $(P,P') \neq 1$. Karena $P$ tak tereduksi dan $\deg P' < \deg P$, satu-satunya kemungkinan ialah $P'=0$.
	
	(iii) $\implies$ (iv): Misalkan $P = \sum_{h \geq 0} b_h X^h$, kondisi $P'=0$ setara dengan $\forall h \in \Z_{\geq 1},\; h b_h = 0$. Karena $\deg P > 0$, kondisi hanya mungkin terpenuhi ketika $p := \text{char}(F) > 0$, pada saat itu $p \nmid h \implies b_h = 0$. % Oleh karena itu $P$ berbentuk $\sum_{k \geq 0} a_k X^{pk}$, di mana $a_k := b_{pk}$.
	
	(iv) $\implies$ (i): Misalkan $P = \sum_{k \geq 0} a_k X^{pk} = P_1(X^p)$, dengan $P_1 := \sum_{k \geq 0} a_k X^k$. Misalkan $y_1, \ldots, y_k \in $ adalah akar-akar $P_1$ dalam $\overline{F}$; untuk setiap $y_i$, ambil $x_i \in \overline{F}$ yang memenuhi $x_i^p = y_i$. Karena $\text{char}(F)=p$, penerapan \eqref{eqn:freshmen-dream} memberikan
	\[ (X^p-y_i) =  (X-x_i)^p, \]
	Jadi, sebagai akar $P$, $x_i$ mempunyai multiplisitas $\geq p > 0$.
\end{proof}

\begin{definition}[Unsur separabel]
	Dalam ekstensi medan $E|F$, jika polinomial minimal $P_x \in F[X]$ dari unsur aljabar $x \in E$ tidak mempunyai akar ganda dalam medan pemisahnya, maka $x$ disebut \emph{separabel} atas $F$.
\end{definition}

Teorema di atas memberikan kriteria awal untuk separabilitas. Untuk analisis lebih lanjut, diasumsikan $\text{char}(F)=p > 0$ dan diingat konstruksi dari teorema \ref{prop:separable-polynomial} (iv): Misalkan $P \in F[X]$ tak tereduksi dan $P'=0$. Maka harus dapat ditulis dalam bentuk $P(X) = P_1(X^p)$. Diperhatikan bahwa
\begin{align*}
	F[X] & \longrightarrow F[X] \\
	f(X) & \longmapsto f(X^p)
\end{align*}
merupakan homomorfisme $F$-aljabar, dan $\deg f(X^p) = p \deg f(X)$. Ketaktereduksian $P$ menjamin bahwa $P_1$ juga tak tereduksi. Jika $P'_1 = 0$, ulangi operasi yang sama untuk memperoleh $P_1(X) = P_2(X^p)$ dengan $P_2$ tetap tak tereduksi. Berdasarkan derajat, iterasi ini harus berhenti sesudah berhingga langkah dan memberikan
\begin{gather}\label{eqn:P-to-P-flat}
	P(X) = P^\flat\left( X^{p^m} \right), \quad P^\flat \in F[X]:\; \text{tak tereduksi}, \; (P^\flat)' \neq 0.
\end{gather}
Uraikan $P^\flat$ di $\overline{F}[X]$ menjadi faktor-faktor linear $c \prod_{i=1}^k (X - y_i)$. Teorema \ref{prop:separable-polynomial} menyiratkan $i = j \iff y_i=y_j$. Untuk setiap $y_i$ ambil $x_i \in \overline{F}$ sehingga $x_i^{p^m} = y_i$. Demikian pula, dari \eqref{eqn:freshmen-dream} diketahui $X^{p^m} - y_i = (X - x_i)^{p^m}$, sehingga $x_i$ unik. Kesimpulan: $P$ di $\overline{F}[X]$ terurai menjadi
\begin{gather*}
	P = c \prod_{i=1}^k (X - x_i)^{p^m}, \quad i=j \iff x_i = x_j
\end{gather*}
Khususnya, setiap akar $P$ dalam medan pemisahnya mempunyai multiplisitas $p^m$. Gagasannya ialah membagi masalah menjadi dua bagian: polinomial $P^\flat$ tanpa akar ganda, atau ``separabel'', dan polinomial ``murni tak separabel'' berbentuk $X^{p^m} - y_i$. Konstruksi ini akan berulang kali digunakan dalam pembahasan separabilitas. Sekarang kita dapat mendefinisikan ekstensi separabel.

\begin{definition}\index[sym1]{$[E:F]_s$}\index{yukuozhang! Derajat separabel (separable degree)}
	Misalkan $E|F$ ekstensi aljabar. \emph{Derajat separabel}-nya didefinisikan sebagai $[E:F]_s := |\Hom_F(E, \overline{F})|$.
\end{definition}
Karena penutup aljabar $F$ unik hingga isomorfisme, definisi $[E:F]_s$ tidak bergantung pada pemilihan $\overline{F}|F$. Selain itu, Korolari \ref{prop:alg-ext-embedding} juga menyiratkan $[E:F]_s \geq 1$.

\begin{proposition}\label{prop:field-tower-sdegree}
	Jika $L|E$ dan $E|F$ merupakan ekstensi aljabar, maka $[L:F]_s = [L:E]_s [E:F]_s$, dengan $[L:E]_s = \Hom_E(L, \overline{F})$ dapat didefinisikan melalui sembarang $F$-pembenaman $\sigma: E \to \overline{F}$; perkaliannya adalah perkalian kardinal.
\end{proposition}
\begin{proof}
	Memperluas pembenaman terpilih $F \hookrightarrow \overline{F}$ ke $\tau: L \to \overline{F}$ dapat dibagi menjadi dua langkah. Langkah pertama adalah memperluas ke $\sigma: E \to \overline{F}$, langkah kedua adalah memperluas setiap $\sigma$ menjadi $\tau: L \to \overline{F}$. Menurut definisi, langkah pertama memiliki tepat $[E:F]_s$ pilihan.

	Selanjutnya pertimbangkan $\sigma \in \Hom_F(E, \overline{F})$ yang diberikan. Karena $F \hookrightarrow \overline{F}$ memberikan ekstensi aljabar, $E \xrightarrow{\sigma} \overline{F}$ juga demikian, dan $\overline{F}$ juga tertutup secara aljabar, maka $E \xrightarrow{\sigma} \overline{F}$ memberikan penutup aljabar dari $E$. Dengan demikian, langkah kedua menurut definisi untuk setiap $\sigma$ ada tepat $[L:E]_s$ pilihan. Mengalikan kardinalitas menghasilkan $[L:F]_s = [L:E]_s [E:F]_s$.
\end{proof}

Derajat separabel ekstensi sederhana $F(x)|F$ dapat dihitung langsung dari polinomial minimal $x$.
\begin{lemma}\label{prop:field-sdegree-F(x)}
	Pertimbangkan ekstensi berhingga berbentuk $F(x)|F$, dan tuliskan polinomial minimal dari $x$ sebagai $P_x$. Maka $[F(x):F]_s$ sama dengan jumlah akar $P_x$ dalam $\overline{F}$ (tanpa menghitung multiplikitas). Ketika $\mathrm{char}(F)=0$ maka $[F(x):F]_s = [F(x):F] = \deg P_x$. Jika $p := \mathrm{char}(F) > 0$, tulis $P_x$ menurut \eqref{eqn:P-to-P-flat}
	\[ P_x(X) = P^\flat_x\left( X^{p^m} \right), \]
	Maka $[F(x):F]_s = \deg P^\flat_x$, $[F(x):F] = p^m [F(x):F]_s$. Secara khusus, $x$ separabel jika dan hanya jika $[F(x):F]_s = [F(x):F]$.
\end{lemma}
\begin{proof}
	Proposisi \ref{prop:field-embedding} menyatakan bahwa $[F(x):F]_s$ sama dengan jumlah akar $P_x$ dalam $\overline{F}$ (tanpa menghitung multiplikitas). Teorema \ref{prop:separable-polynomial} menunjukkan bahwa akar ganda hanya mungkin ketika $p := \mathrm{char}(F) > 0$. Dalam pembahasan \eqref{eqn:P-to-P-flat} telah ditunjukkan bahwa terdapat tepat $\deg P^\flat_x$ akar berbeda, masing-masing bermultiplikitas $p^m$.
\end{proof}

Untuk ekstensi berhingga umum $E|F$, Korolari \ref{prop:alg-ext-embedding} menyatakan $[E:F]_s \leq [E:F]$; kita bahkan memperoleh sifat keterbagian berikut.
\begin{definition-theorem}\label{def:insep-deg}
	Berikut ini hanya mempertimbangkan ekstensi berhingga.
	\begin{enumerate}[(i)]
		\item Untuk setiap $E|F$, berlaku $[E:F]_s \mid [E:F]$. Karena itu \emph{derajat tak separabel} $E|F$ dapat didefinisikan sebagai \index{yukuozhang!derajat tak separabel (inseparable degree)}\index[sym1]{$[E:F]_i$}
			\[ [E:F]_i := \frac{[E:F]}{[E:F]_s}. \]
		\item Hanya mungkin ada $[E:F]_i > 1$ saat $p := \text{char}(F) > 0$, dan pada saat itu itu pasti dalam bentuk $p^m$.
		\item Derajat tak separabel ekstensi berhingga memenuhi $[L:F]_i = [L:E]_i [E:F]_i$.
		\item Berikut setara:
			\begin{inparaenum}[(a)]
				\item $E$ mempunyai suatu keluarga pembangkit yang separabel atas $F$;
				\item Persamaan $[E:F]_s = [E:F]$ berlaku,
				\item setiap unsur $E$ separabel atas $F$.
			\end{inparaenum}
		\item Untuk ekstensi sederhana $F(x)|F$, multiplisitas setiap akar polinomial minimal $P_x$ dalam $\overline{F}$ sama dengan $[F(x):F]_i$; khususnya, $x$ separabel atas $F$ jika dan hanya jika $[F(x):F]_i = 1$.
	\end{enumerate}
\end{definition-theorem}
\begin{proof}
	Kita buktikan (i) terlebih dahulu. Misalkan $E = F(x_1, \ldots, x_n)$ dan tetapkan $F_m := F(x_1, \ldots, x_m)$. Dengan membandingkan Proposisi \ref{prop:field-tower-degree} dan \ref{prop:field-tower-sdegree}, cukup dibuktikan $[F_m:F_{m-1}]_s | [F_m:F_{m-1}]$ untuk setiap ekstensi antara $F_m|F_{m-1}$. Karena itu, kita boleh mengandaikan $E = F(x)$; kasus ini telah ditangani oleh Lemma \ref{prop:field-sdegree-F(x)}, yang sekaligus memberikan (ii).
	
	Properti (iii) adalah kesimpulan langsung dari properti yang sesuai dari $[L:F]$ dan $[L:F]_s$.

	Argumen di atas menunjukkan: jika semua pembangkit $E$, yaitu $x_m$, separabel atas $F$, maka tentu juga separabel atas $F_{m-1}$, sehingga $[E:F] = \prod_{m=1}^n [F_m : F_{m-1}] = \prod_{m=1}^n [F_m : F_{m-1}]_s = [E:F]_s$. Jika diasumsikan $[E:F]=[E:F]_s$, maka untuk sembarang $x \in E$ berlaku $[F(x):F]_i \mid [E:F]_i = 1$, dan Lemma \ref{prop:field-sdegree-F(x)} memastikan bahwa $x$ separabel. Dengan ini (iv) terbukti. Akhirnya, (v) merupakan akibat langsung Lemma \ref{prop:field-sdegree-F(x)}.
\end{proof}

\begin{definition}\index{yukuozhang! separabel (separable)}
	Jika setiap unsur dalam ekstensi aljabar $E|F$ separabel atas $F$, maka disebut \emph{ekstensi separabel}. Jadi ekstensi berhingga $E|F$ separabel jika dan hanya jika $[E:F]_s = [E:F]$, atau secara ekuivalen $[E:F]_i = 1$.
\end{definition}

\begin{proposition}\label{prop:sep-ext-dist}
	Ekstensi separabel bersifat istimewa dalam arti Konvensi \ref{con:dist-extensions} dan tertutup terhadap kompositum.
\end{proposition}
\begin{proof}
	Untuk membuktikan sifat \textbf{D.1}, lakukan pengamatan sederhana berikut:
	\begin{compactitem}
		\item Jika $x \in L$ separabel atas $F$ (yakni polinomial minimalnya tidak mempunyai akar ganda dalam $\overline{F}$), maka tentu ia juga separabel atas $E$. Jadi, $L|F$ separabel menyiratkan bahwa $L|E$ dan, jelas pula, $E|F$ separabel.
		\item Jika $L|E$ dan $E|F$ separabel, untuk sembarang $x \in L$ perhatikan polinomial minimal $Q_x \in E[X]$. Koefisien-koefisiennya membangkitkan subekstensi berhingga $E_0|F$ dari $E|F$. Bagian (iv) Definisi--Teorema \ref{def:insep-deg} menunjukkan bahwa $E_0|F$ separabel, sedangkan $x$ separabel atas $E_0$. Karena $E_0(x) \supset F(x)$, bagian (iii) dan (v) Definisi--Teorema \ref{def:insep-deg} memberikan
			\[ [F(x):F]_i \quad \text{membagi} \quad [E_0(x):F]_i = [E_0(x):E_0]_i [E_0:F]_i = 1, \]
			Oleh karena itu, $x$ separabel atas $F$.
	\end{compactitem}
	Sifat \textbf{D.2}, \textbf{D.3}, dan ketertutupan terhadap kompositum mengikuti sifat berikut: jika ekstensi aljabar $E|F$ mempunyai suatu keluarga pembangkit yang separabel atas $F$, maka $E|F$ separabel. Memang, setiap $x$ dapat ditulis menggunakan berhingga banyak pembangkit, sehingga masalah direduksi ke kasus ekstensi berhingga dan bagian (iv) Definisi--Teorema \ref{def:insep-deg} dapat diterapkan.
\end{proof}

Ketertutupan ekstensi separabel terhadap kompositum memungkinkan kita mendefinisikan suatu penutup.
\begin{definition}\label{def:sep-closure}\index{kefenbibao@penutup separabel (separable closure)}
	Kompositum semua subekstensi separabel dari suatu ekstensi $\Omega|F$ tetap separabel dan disebut penutup separabel $F$ di dalam $\Omega$. Penutup separabel di dalam penutup aljabar $\overline{F}|F$ disebut singkatnya \emph{penutup separabel} dari $F$ dan ditulis $F^\mathrm{sep}|F$.
\end{definition}
Penutup separabel juga mempunyai karakterisasi intrinsik berikut, tanpa perlu melalui $\overline{F}|F$.
\begin{proposition}
	Penutup separabel $F^\mathrm{sep}|F$ adalah ekstensi normal: tepatnya, ia merupakan medan pemisah dari $\mathcal{P} := \{P \in F[X]: \text{tak tereduksi},\; P' \neq 0 \}$.
\end{proposition}
Selain itu, juga dapat meniru cara mendefinisikan penutup aljabar, mendefinisikan $F^\text{sep}|F$ sebagai ekstensi separabel maksimum dari $F$; lihat Lemma \ref{prop:alg-closed-no-ext}.

\begin{definition}\label{def:perfect-field}\index{yu! Sempurna}
	Jika semua ekstensi aljabar dari medan $F$ separabel, maka $F$ disebut \emph{medan sempurna}.
\end{definition}

Untuk studi mengenai kasus $p := \text{char}(F) > 0$, \eqref{eqn:freshmen-dream} akan memainkan peran kunci; ia memastikan sifat-sifat berikut.
\begin{itemize}
	\item Untuk setiap $m \geq 0$, definisikan $F^{p^m} := \left\{x^{p^m} : x \in F \right\}$. Ini adalah submedan dari $F$: memang, $F^{p^m}$ tertutup terhadap perkalian dan pengambilan invers unsur tak nol $y \mapsto y^{-1}$, memuat submedan prima $\F_p$, serta tertutup terhadap penjumlahan dan invers aditif $y \mapsto -y = \underbracket{(-1)}_{\in \F_p} \cdot y$.
	\item Setiap $a \in F$ di $\overline{F}$ memiliki tepat satu akar $p^m$, sebut saja $a^{p^{-m}}$.
\end{itemize}

\begin{lemma}\label{prop:pins-polynomial}
	Misalkan $p := \text{char}(F) > 0$, $m \geq 1$. Polinomial $X^{p^m} - a \in F[X]$ tak tereduksi jika dan hanya jika $a \notin F^p$.
\end{lemma}
\begin{proof}
	Jika $a = b^p$ dengan $b \in F$, maka $X^{p^m}-a = \left( X^{p^{m-1}}-b \right)^p$ tereduksi. Berikutnya andaikan $a \notin F^p$.

	Misalkan $P$ polinomial minimal dari $\alpha := a^{p^{-m}} \in \overline{F}$ atas $F$. Karena setiap faktor prima dari $X^{p^m}-a =(X-\alpha)^{p^m}$ dalam $F[X]$ mempunyai akar $\alpha$, semuanya habis dibagi oleh $P$. Faktorisasi unik dalam $F[X]$ memberikan bilangan bulat positif $k$ dengan $X^{p^m}-a = P^k$; kita harus menunjukkan $k=1$. Perbandingan derajat memberikan $k = p^n$ ($n \leq m$), lalu perbandingan suku konstan memberikan $-a \in F^{p^n}$, sehingga $a \in F^{p^n}$. Asumsi memaksa $n=0$. Terbukti.
\end{proof}
Khususnya, jika $x \in E \smallsetminus F$ dan $a := x^p \in F$, Lemma \ref{prop:pins-polynomial} menyiratkan bahwa $X^p - a \in F[X]$ adalah polinomial minimal $x$. Lemma \ref{prop:field-sdegree-F(x)} langsung memberikan $[F(x):F]_s = 1$. Dalam \S\ref{sec:pins} kita akan membahas polinomial semacam ini lebih lengkap. Sekarang kembali ke medan sempurna.

\begin{theorem}\label{prop:perfect-field-char}\index{daishu!sempurna}
	Medan $F$ sempurna jika dan hanya jika
	\begin{inparaenum}[(a)]
		\item $\text{char}(F)=0$, atau
		\item $p := \text{char}(F) > 0$ dan $F = F^p$.
	\end{inparaenum}
\end{theorem}
\begin{proof}
	Berdasarkan teorema \ref{prop:separable-polynomial}, hanya perlu memverifikasi $F=F^p$ dalam kasus $p := \text{char}(F) > 0$ sama dengan semua $P \in F[X]$ yang tak tereduksi memenuhi $P' \neq 0$. Pertama, misalkan $F=F^p$, jika $P' = 0$ maka terdapat ekspresi $P = \sum_{k \geq 0} a_k X^{pk}$, ambil $a'_k \in F$ sehingga $(a'_k)^p = a_k$, ini akan mengakibatkan $P = (\sum_{k \geq 0} a'_k X^k)^p$ tereduksi.
	
	Sebaliknya, misalkan $\exists a \in F \smallsetminus F^p$. Lemma \ref{prop:pins-polynomial} menunjukkan bahwa $P := X^p - a \in F[X]$ tak tereduksi, sekaligus $P' = pX^{p-1} = 0$. Terbukti.
\end{proof}

\begin{corollary}
	Semua medan berhingga adalah medan sempurna.
\end{corollary}
\begin{proof}
	Karakteristik medan berhingga $F$ harus suatu bilangan prima $p$, sebab jika tidak maka $\Q \hookrightarrow F$. Peta $x \mapsto x^p$ adalah endomorfisme grup aditif $F$ dengan kernel $\{x \in F: x^p=0 \} = \{0\}$, sehingga injektif. Sifat berhingga kemudian memberikan $F=F^p$.
\end{proof}

\section{Teorema Unsur Primitif}
Di awal bagian ini, kita akan terlebih dahulu meninjau salah satu hasil dari latihan di Bab Lima.
\begin{theorem}\label{prop:field-subgroup-cyclic}
	Misalkan $F$ medan. Setiap subgrup berhingga dari $F^\times$ adalah siklik. Khususnya, jika $F$ medan berhingga, maka $F^\times$ siklik.
\end{theorem}
\begin{proof}
	Dalam latihan \S\ref{sec:ring} telah digambarkan pendekatan berdasarkan polinomial siklotom; di sini kita memberikan bukti lain menggunakan teorema struktur grup abelian yang dibangkitkan secara berhingga (Korolari \ref{prop:finite-abelian-structure}). Misalkan $A \subset F^\times$ grup berhingga. Menurut hasil tersebut, terdapat bilangan bulat $> 1$, yaitu $d_1 \mid \cdots \mid d_n$, sedemikian sehingga $A \simeq \bigoplus_{i=1}^n \Z/d_i\Z$. Maka $\left\{a\in A: a^{d_n}=1 \right\}$ mempunyai $\prod_{i=1}^n d_i$ unsur. Namun polinomial $X^{d_n}-1$ mempunyai paling banyak $d_n$ akar dalam $F$, sehingga $n=1$ dan $A$ memang siklik.
\end{proof}

\begin{definition}\index{benyuanyuansu@unsur primitif (primitive element)}
	Misalkan $E|F$ ekstensi berhingga. Jika $u \in E$ memenuhi $E=F(u)$, maka $u$ disebut \emph{unsur primitif} dari $E|F$.
\end{definition}

\begin{example}
	Pertama coba kasus medan bilangan kuadrat. Misalkan $p,q$ bilangan prima yang berbeda. Mudah dilihat bahwa $p$ tidak mempunyai akar kuadrat dalam $\Q(\sqrt{q}) = \Q + \Q\sqrt{q}$, sehingga $[\Q(\sqrt{p},\sqrt{q}): \Q] = [\Q(\sqrt{p},\sqrt{q}):\Q(\sqrt{q})] [\Q(\sqrt{q}):\Q] = 4$. Verifikasi bahwa $u := \sqrt{p}+\sqrt{q}$ merupakan unsur primitif dari $\Q(\sqrt{p}, \sqrt{q})|\Q$ sebagai berikut:
	\[ u^3 = \left( \sqrt{p}+\sqrt{q} \right)^3 = p\sqrt{p} + 3p\sqrt{q} + 3q\sqrt{p} + q\sqrt{q} = (p+3q)\sqrt{p} + (3p+q)\sqrt{q}, \]
	Karena itu $\sqrt{q} = (2p-2q)^{-1}(u^3 - (p+3q)u) \in \Q(u)$; dari $\sqrt{p} = u - \sqrt{q} \in \Q(u)$ diketahui bahwa $u$ memang unsur primitif.
\end{example}
Untuk ekstensi medan yang lebih umum, perhitungan serupa segera menjadi sangat rumit. Untungnya, kita mempunyai teorema umum berikut.

\begin{theorem}\label{prop:prim-element-separable}
	Ekstensi separabel berhingga $E|F$ pasti memiliki unsur primitif. Ketika $F$ tak terhingga dan $E=F(x_1, \ldots, x_n)$, hasilnya dapat diperkuat lebih lanjut menjadi: Misalkan $x_2, \ldots, x_n$ separabel atas $F$, maka unsur primitif dapat diambil sebagai kombinasi linear $x_1, \ldots, x_n$ atas $F$.
\end{theorem}
\begin{proof}
	Jika $F$ berhingga maka $E$ juga berhingga. Teorema \ref{prop:field-subgroup-cyclic} memberikan pembangkit $u$ dari $E^\times$; jelas $E^\times = \lrangle{u} \implies E=F(u)$. Kasus medan berhingga terbukti.
	
	Misalkan $F$ tak hingga dan $E=F(x_1, \ldots, x_n)$; pertama-tama periksa kasus $n=2$. Tuliskan $P_1, P_2$ sebagai polinomial minimal masing-masing $x_1, x_2$. Sekarang akan dibuktikan bahwa untuk $t \in F$ yang ``umum'' berlaku $E=F(x_1 + tx_2)$; perhatikan bahwa $x_1 = (x_1 + tx_2) - tx_2$, jadi cukup membuktikan $x_2 \in F(x_1 + tx_2)$.
	
	Membenamkan $F(x_1, x_2)$ ke dalam penutup aljabar $\overline{F}$ dan mempertimbangkan unsur di $F(x_1 + tx_2)[X]$
	\[ P_1(x_1 + tx_2 - tX), \quad P_2(X). \]
	Karena keduanya memiliki akar bersama $x_2$ dalam $\overline{F}$, lema \ref{prop:coprime-basechange} menyiratkan bahwa faktor persekutuan utama terbesar $R \in F(x_1 + tx_2)[X]$ memenuhi $\deg R \geq 1$. Jika dapat mengambil $t$ sedemikian sehingga $\deg R=1$ maka pasti ada $R = X-x_2$, pada saat itu dapat dihasilkan $x_2 \in F(x_1 + tx_2)$ seperti yang diinginkan.
	
	Asumsikan kebalikan $\deg R > 1$. Dari $R \mid P_2$ dan fakta bahwa $P_2$ tidak memiliki akar ganda (di sini digunakan separabilitas), diketahui bahwa $R$ mempunyai, di dalam $\overline{F}$, suatu akar $y \neq x_2$. Demikian pula, $R \mid P_1(x_1 + tx_2 - tX)$ menyiratkan $P_1(x_1 + t(x_2-y))=0$, sehingga diperoleh
	\begin{gather*}
		x_1 + t(x_2 - y) = z, \quad z \in \overline{F}:\; P_1(z)=0.
	\end{gather*}
	Untuk setiap akar $P_1$ dari $z$ dan akar $P_2$ dari $y \neq x_2$, persamaan di atas menentukan secara unik $t$. Karena jumlah akar berhingga sedangkan $F$ tak terhingga, selalu mungkin memilih $t$ sehingga persamaan di atas tidak berlaku untuk setiap pasangan $(z,y)$, maka $\deg R = 1$.
	
	Ketika $n > 2$, dapat menggunakan $F(x_1, \ldots, x_n) = F(x_1, \ldots, x_{n-1})(x_n)$ untuk membuktikan secara rekursif.
\end{proof}

Karakterisasi berikut memberikan syarat perlu dan cukup bagi keberadaan unsur primitif. Agar mudah dinyatakan, berikutnya andaikan $F \subset E$.
\begin{theorem}[E.\ Steinitz]\label{prop:prim-element}
	Ekstensi berhingga $E|F$ merupakan ekstensi sederhana (yakni mempunyai unsur primitif) jika dan hanya jika hanya ada berhingga banyak medan antara $F \subset M \subset E$.
\end{theorem}
\begin{proof}
	Jika $F$ berhingga, maka $E$ juga berhingga, sehingga tentu hanya ada berhingga banyak medan antara; Teorema \ref{prop:prim-element-separable} menyatakan keberadaan unsur primitif. Berikutnya cukup ditangani kasus $F$ tak hingga.
	
	Pertama, andaikan $E = F(u)$. Untuk setiap medan antara $M$, tuliskan polinomial minimal $u$ atas $M$ sebagai $P_M \in M[X]$; maka $P_M \mid P_F$. Kita nyatakan bahwa $M|F$ dibangkitkan oleh koefisien-koefisien $c_0, \ldots, c_m$ dari $P_M$. Memang, karena $P_M$ tak tereduksi atas $M$, polinomial itu juga tak tereduksi atas $F(c_0, \ldots, c_m) \subset M$, sehingga
	\[ [E:M] = \deg P_M = [E:F(c_0, \ldots, c_m)], \]
	Sehingga $M=F(c_0, \ldots, c_m)$, yang juga $M$ dapat direkonstruksi dari $P_M$. Dari sini terlihat pemetaan
	\[\begin{tikzcd}[row sep=tiny, column sep=small]
		\{ \text{medan antara}\; F \subset M \subset L \} \arrow[r] & \{ P_F \;\text{adalah faktor monik dalam $E[X]$} \} \\
		M \arrow[mapsto, r] & P_M
	\end{tikzcd}\]
	bersifat injektif. Karena sisi kanan berhingga, sisi kiri juga berhingga.
	
	Sekarang andaikan hanya ada berhingga banyak medan antara $M$. Setiap $M$ juga merupakan subruang vektor-$F$ dari $E$. Karena $F$ tak hingga, kita nyatakan bahwa berhingga banyak subruang sejati tidak dapat menutupi seluruh $E$. Maka setiap $u \in E \smallsetminus \bigcup\{M : M \subsetneq E \}$ merupakan unsur primitif. Untuk membuktikan pernyataan itu, andaikan sebaliknya. Maka $E$ dapat ditutupi oleh berhingga banyak hiperbidang linear, yakni
	\[ E = \bigcup_{i=1}^k \{x \in E : L_i(x) = 0\}, \]
	Di antara ini, $L_1, \ldots, L_k: E \to F$ adalah beberapa pemetaan $F$-linear yang tidak nol; maka $\prod_{i=1}^k L_i$ adalah fungsi polinomial derajat $k$ yang selalu bernilai nol pada $E \simeq F^{[E:F]}$, yang bertentangan dengan proposisi \ref{prop:polynomial-function}.
\end{proof}
Dengan menggunakan terminologi geometri aljabar, bukti di paragraf terakhir sebenarnya menunjukkan bahwa dalam ekstensi tunggal $E|F$, unsur 'umum' $E$ (yaitu: kecuali sejumlah berhingga $M \subsetneq E$) semuanya adalah unsur primitif.

Dalam latihan akan diberikan contoh konkret ekstensi tanpa unsur primitif, dan akan diverifikasi bahwa dalam kasus itu memang terdapat tak hingga banyak medan antara.

\section{Norma dan Teras dalam Ekstensi Medan}\label{sec:norm-trace-fields} \index{ji-trace}\index{fanshu}
Untuk ekstensi berhingga $E|F$, Definisi \ref{def:norm-trace} mendefinisikan, bagi setiap $F$-aljabar $E$, dua peta dari $E$ ke $F$, yakni norma $\Nm_{E|F}$ dan teras $\Tr_{E|F}$; masing-masing merupakan homomorfisme monoid perkalian dan grup aditif. Untuk menara ekstensi medan $L|E|F$, Teorema \ref{prop:norm-trace-transitivity} memberikan
\[ \Nm_{L|F} = \Nm_{E|F} \Nm_{L|E}, \quad \Tr_{L|F}=\Tr_{E|F}\Tr_{L|E}. \]

Untuk ekstensi medan $E|F$ dan unsur $x \in E$, perhitungan norma dan terasnya dapat dibagi menjadi dua tahap:
\[ \Nm_{E|F}(x) = \Nm_{F(x)|F}(\Nm_{E|F(x)}(x)), \quad \Tr_{E|F}(x)=\Tr_{F(x)|F}(\Tr_{E|F(x)}(x)), \]
Adapun tahap $E|F(x)$ mudah: karena $x \in F(x)$, sifat-sifat yang tercantum dalam \S\ref{sec:trace-norm-disc} menunjukkan
\begin{gather}\label{eqn:norm-trace-field-prepa}
	\Nm_{E|F(x)}(x) = x^{[E:F(x)]}, \quad \Tr_{E|F(x)}(x) = [E:F(x)] x.
\end{gather}
Seperti biasa, strategi yang masuk akal adalah memulai dengan memeriksa $\Nm_{F(x)|F}(x)$ dan $\Tr_{F(x)|F}(x)$.

\begin{theorem}\label{prop:norm-trace-F(x)}
	Misalkan $x$ memiliki polinomial minimal $P_x = X^n + a_{n-1}X^{n-1} + \cdots + a_0$ di atas $F$, maka
	\[ \Nm_{F(x)|F}(x) = (-1)^n a_0, \quad \Tr_{F(x)|F}(x) = -a_{n-1}. \]
\end{theorem}
\begin{proof}
	Gunakan Definisi \ref{def:norm-trace} semula. Ambil basis berurutan $1, x, \ldots, x^{n-1}$ dari ruang vektor-$F$ $F(x)$. Karena $x \cdot x^{n-1} = \sum_{i=0}^{n-1} (-a_i)x^i$, matriks peta perkalian kiri $m_x: a \mapsto xa$ adalah
	\[ A := \begin{pmatrix}
		0 & \cdots & 0 & -a_0 \\
		1 & & & -a_1 \\
		& \ddots & & \vdots \\
		& & 1 & -a_{n-1}
	\end{pmatrix}\]
	Mudah dihitung $\Nm_{F(x)|F}(x) = \det(A) = (-1)^n a_0$, sedangkan $\Tr_{F(x)|F}(x) = \Tr(A) = -a_{n-1}$.
\end{proof}

\begin{theorem}\label{prop:norm-trace-field}
	Untuk ekstensi berhingga $E|F$ dan $x \in E$, ambil penutup aljabar $\overline{F}|F$, maka ada
	\begin{align*}
		\Nm_{E|F}(x) & = \prod_{\sigma \in \Hom_F(E, \overline{F})} \sigma(x)^{[E:F]_i}, \\
		\Tr_{E|F}(x) & = [E:F]_i \sum_{\sigma \in \Hom_F(E, \overline{F})} \sigma(x) \\
		& = \begin{cases}
			\sum_{\sigma \in \Hom_F(E, \overline{F})} \sigma(x), & E|F\; \text{separabel} \\
			0, & E|F\; \text{tak separabel}.
		\end{cases}
	\end{align*}
\end{theorem}
\begin{proof}
	Tuliskan polinomial minimal $x$ sebagai $P_x = X^n + a_{n-1}X^{n-1} + \cdots + a_0$. Dalam Teorema \ref{prop:norm-trace-F(x)}, $(-1)^n a_0$ dan $-a_{n-1}$ masing-masing adalah hasil kali dan jumlah semua akar $P_x$ dalam $\overline{F}$, dengan memperhitungkan multiplikitas; multiplikitas tiap akar sama dengan $[F(x):F]_i$ (Definisi--Teorema \ref{def:insep-deg} (v)). Dengan Proposisi \ref{prop:field-embedding}, ini dapat ditulis ulang sebagai
	\begin{align*}
		\Nm_{F(x)|F}(x) & = \prod_{\sigma \in \Hom_F(F(x), \overline{F})} \sigma(x)^{[F(x):F]_i} , \\
		\Tr_{F(x)|F}(x) & = [F(x):F]_i \sum_{\sigma \in \Hom_F(F(x), \overline{F})} \sigma(x).
	\end{align*}
	Setiap $\sigma \in \Hom_F(F(x), \overline{F})$ ke $E$ mempunyai sebanyak $[E: F(x)]_s$ perluasan. Karena itu \eqref{eqn:norm-trace-field-prepa} memberikan
	\begin{align*}
		\Nm_{E|F}(x) & = \Nm_{F(x)|F}(x)^{[E:F(x)]} = \prod_{\sigma \in \Hom_F(F(x), \overline{F})} \sigma(x)^{[F(x):F]_i [E:F(x)]} \\
		& = \prod_{\sigma \in \Hom_F(E, \overline{F})} \sigma(x)^{[F(x):F]_i [E:F(x)]_i [E:F(x)]_s \big/ [E:F(x)]_s} \\
		& = \prod_{\sigma \in \Hom_F(E, \overline{F})} \sigma(x)^{[E:F]_i}.
	\end{align*}
	Dalam operasi yang sama, mengganti perkalian dengan penjumlahan akan menghasilkan $\Tr_{E|F}(x) = [E:F]_i \sum_{\sigma \in \Hom_F(E, \overline{F})} \sigma(x)$. Karena $E|F$ separabel jika dan hanya jika $[E:F]_i = 1$, sedangkan kondisi tak separabel pasti memiliki karakteristik $p > 0$ serta $[E:F]_i = p^m$, $m \geq 1$, pada saat ini $[E:F]_i \cdot 1_{\overline{F}} = 0$.
\end{proof}

Rumus-rumus di atas membantu menjelaskan bentuk teras yang dibahas sesudah Definisi \ref{def:discriminant}
\begin{align*}
	\Tr_{E|F}: E \times E & \longrightarrow F \\
	(x,y) & \longmapsto \Tr_{E|F}(xy).
\end{align*}
Ini adalah bentuk bilinear simetris pada $E$. Secara lebih umum, misalkan $V$, $W$ ruang vektor berdimensi berhingga atas $F$ dan $B: V \times W \to F$ peta bilinear. Jika
\[ B(v,\cdot) = 0 \iff v=0, \quad B(\cdot,w)=0 \iff w=0 \]
maka $B$ disebut tak terdegenerasi. Sekarang andaikan $\dim_F V = \dim_F W = n$. Jika basis $x_1, \ldots, x_n$ dari $V$ dan basis $y_1, \ldots, y_n$ dari $W$ (dengan urutan diperhitungkan) memenuhi
\[ B(x_i, y_j) = \begin{cases} 1, & i=j \\ 0, & i \neq j \end{cases} \]
maka kedua basis disebut saling dual. Teori dasar aljabar linear menyiratkan bahwa $V$, $W$ mempunyai sepasang basis dual jika dan hanya jika $B$ tak terdegenerasi; dalam hal ini setiap $(x_i)_i$ mempunyai basis dual yang unik.

\begin{theorem}
	Misalkan $E=F(x)$, dengan polinomial minimal $P_x$ dari $x$ berderajat $n$; tuliskan
	\[ \frac{P_x}{X-x} = \sum_{i=0}^{n-1} b_i X^i \quad \in E[X]. \]
	Jika $P'_x(x) \neq 0$, maka untuk basis $1, x, \ldots, x^{n-1}$ dari $E$, basis dual terhadap $\Tr_{E|F}$ adalah
	\[ \frac{b_0}{P'_x(x)}, \ldots, \frac{b_{n-1}}{P'_x(x)}. \]
\end{theorem}
\begin{proof}
	Ambil penutup aljabar $\overline{F}$ dan benamkan $E$ ke dalamnya. Misalkan $x_1, \ldots, x_n \in \overline{F}$ adalah akar-akar $P_x$. Teorema \ref{prop:separable-polynomial} memastikan bahwa akar-akar itu berbeda dan bahwa $E|F$ separabel. Perhatikan pula bahwa untuk setiap $1 \leq i,k \leq n$ berlaku
	\[ \left. \frac{P_x}{X - x_i} \right|_{X=x_k} = \begin{cases}
		P'_x(x_k), & k =i \\
		0, & k \neq i.
	\end{cases}\]
	(Ini elementer; lihat \eqref{eqn:derivation-linear-approx}). Pertama-tama kita nyatakan
	\[ \sum_{i=1}^n \frac{P_x}{X-x_i} \cdot \frac{x_i^j}{P'_x(x_i)} = X^j, \quad 0 \leq j \leq n-1. \]
	Untuk setiap $j$, dari pengamatan di atas diketahui bahwa selisih kedua ruas merupakan polinomial yang mempunyai $n$ akar $x_1, \ldots, x_n$ dan berderajat $\leq n-1$, sehingga kedua ruas sama. Persamaan ini juga dapat ditulis sebagai
	\[ \sum_{\sigma \in \Hom_F(E,\overline{F})} \sigma \left( \frac{P_x}{X-x} \cdot \frac{x^j}{P'_x(x)} \right) = X^j \qquad
	\begin{tikzpicture}[baseline]
			\matrix[matrix of math nodes, left delimiter={\|}]
				{\sigma \leftrightarrow i \\ \sigma(x)=x_i \\ };
	\end{tikzpicture} \]
	Di sini $\sigma$ diterapkan pada koefisien polinomial dan diperpanjang secara alami menjadi homomorfisme $F$-aljabar $E[X] \xrightarrow{\sigma} \overline{F}[X]$. Dengan membandingkan koefisien $X^i$ pada kedua ruas dan menerapkan Teorema \ref{prop:norm-trace-field}, diperoleh
	\[ \Tr_{E|F}\left(  x^j \cdot \frac{b_i}{P'_x(x)}\right)=\sum_{\sigma \in \Hom_F(E,\overline{F})} \sigma \left( b_i \cdot \frac{x^j}{P'_x(x)} \right)  = \begin{cases} 1, & i=j, \\ 0, & i \neq j \end{cases} \]
	Inilah syarat bagi basis dual.
\end{proof}

Untuk ekstensi berhingga separabel $E|F$, Teorema \ref{prop:prim-element-separable} memastikan bahwa prasyarat teorema di atas terpenuhi; untuk ekstensi berhingga tak separabel selalu berlaku $\Tr_{E|F}=0$. Korolari berikut mengalir langsung darinya.
\begin{corollary}\label{prop:separable-discriminant}
	Bentuk teras $\Tr_{E|F}: E \times E \to F$ dari ekstensi berhingga $E|F$ tak terdegenerasi jika dan hanya jika $E|F$ separabel.
\end{corollary}

Jika tujuan kita hanya membuktikan bahwa $\Tr_{E|F}$ tak terdegenerasi, Teorema \ref{prop:indep-character-Artin} yang diperkenalkan kemudian memberikan bukti lebih singkat. Namun, uraian eksplisit basis dual sangat berguna, misalnya dalam teori bilangan aljabar, dan tidak diberikan oleh argumen abstrak tersebut. Kita serahkan pembuktiannya sebagai latihan.

\section{Ekstensi Murni Tak Separabel}\label{sec:pins}
Karena fenomena tak separabel hanya muncul dalam karakteristik $p > 0$, sepanjang bagian ini tetapkan bilangan prima $p$ dan anggap semua medan berkarakteristik $p$.
\begin{definition}[Unsur murni tak separabel]
	Misalkan $\Omega|F$ ekstensi medan. Jika polinomial minimal suatu unsur aljabar $x \in \Omega$ berbentuk $P_x = X^{p^m}-a \in F[X]$, dengan $m \geq 0$, maka $x$ disebut \emph{unsur murni tak separabel} atas $F$.
\end{definition}
Dari Lemma \ref{prop:field-sdegree-F(x)}, $x$ murni tak separabel jika dan hanya jika $[F(x):F] = [F(x):F]_i = p^m$, yang juga setara dengan $[F(x):F]_s = 1$. Dalam hal ini, polinomial separabel yang diberikan oleh \eqref{eqn:P-to-P-flat}, yaitu $P^\flat_x$, adalah polinomial linear $X-a$. Berikut ini akibat langsung Lemma \ref{prop:pins-polynomial}.

\begin{proposition}
	Misalkan ada $k \geq 0$ sehingga $x^{p^k} \in F$. Ambil $m := \min\left\{k \geq 0 : x^{p^k} \in F \right\}$, $a := x^{p^m}$; maka $x$ adalah unsur murni tak separabel dengan polinomial minimal $X^{p^m}-a$.
\end{proposition}

\begin{definition-theorem}\label{def:pins}\index{yukuozhang!murni tak separabel (purely inseparable)}
	Untuk ekstensi aljabar $E|F$, kondisi berikut setara.
	\begin{enumerate}[(i)]
		\item setiap unsur $E$ murni tak separabel atas $F$;
		\item $E|F$ dibangkitkan oleh suatu keluarga unsur murni tak separabel;
		\item $[E:F]_s = 1$ (ingat bahwa derajat separabel didefinisikan untuk setiap ekstensi aljabar);
	\end{enumerate}
	Jika salah satu kondisi ini terpenuhi, $E|F$ disebut \emph{ekstensi murni tak separabel}.
\end{definition-theorem}
\begin{proof}
	(i) $\implies$ (ii): Jelas.

	(ii) $\implies$ (iii): Ambil penutup aljabar $\overline{F}|F$ dan ingat bahwa $\Hom_F(E, \overline{F})$ tak kosong (Korolari \ref{prop:alg-ext-embedding}). Setiap $\iota \in \Hom_F(E, \overline{F})$ ditentukan oleh tindakannya pada suatu keluarga pembangkit. Pembahasan sebelumnya menunjukkan bahwa polinomial minimal setiap pembangkit murni tak separabel $x$ berbentuk $X^{p^m} - a$. Karena itu $\iota \in \Hom_F(E, \overline{F})$ harus memetakan $x$ ke satu-satunya akar $a^{p^{-m}}$ dari $X^{p^m} - a$ dalam $\overline{F}$. Ini berarti $[E:F]_s = |\Hom_F(E, \overline{F})| = 1$.
	
	(iii) $\implies$ (i): Misalkan $x \in E$. Proposisi \ref{prop:field-tower-sdegree} memberikan $1 = [E:F]_s = [E:F(x)]_s [F(x):F]_s$, sehingga $[F(x):F]_s = 1$. Menurut pembahasan sebelumnya, ini setara dengan $x$ murni tak separabel.
\end{proof}

Tetapkan penutup aljabar $\overline{F}|E$ dan definisikan submedan $\overline{F}$ berikut (verifikasi mudah diserahkan kepada pembaca):
\[ F^{p^{-\infty}} := \{ x \in \overline{F}: \exists m \geq 1,\; x^{p^m} \in F \}  \]
Maka $E|F$ murni tak separabel jika dan hanya jika $E \subset F^{p^{-\infty}}$. Jadi $F^{p^{-\infty}}$ adalah medan pemisah semua polinomial atas $F$ yang berbentuk $X^{p^m}-a$; ia juga disebut \emph{penutup sempurna} dari $F$. Latihan akan menjelaskan asal istilah ini.

\begin{lemma}
	Ekstensi murni tak separabel pasti merupakan ekstensi normal.
\end{lemma}
\begin{proof}
	Gunakan karakterisasi ekstensi normal (\textbf{N.1}): polinomial minimal setiap unsur dalam ekstensi murni tak separabel hanya mempunyai satu akar.
\end{proof}

\begin{proposition}\label{prop:pins-dist}
	Ekstensi murni tak separabel bersifat istimewa dalam arti Konvensi \ref{con:dist-extensions} dan tertutup terhadap kompositum.
\end{proposition}
\begin{proof}
	Pertama tetapkan sifat \textbf{D.1}. Ambil penutup aljabar $\overline{F}|L$. Jika $L \subset E^{p^{-\infty}}$, $E \subset F^{p^{-\infty}}$, maka dari definisi jelas $L \subset F^{p^{-\infty}}$. Sebaliknya, jika $L \subset F^{p^{-\infty}} \implies (L \subset E^{p^{-\infty}}) \wedge (E \subset F^{p^{-\infty}})$ maka lebih jelas.

	Untuk \textbf{D.2}, perhatikan bahwa $L|F$ murni tak separabel menyiratkan setiap $x \in L$ murni tak separabel atas $M$ (karena $x^{p^m} \in F \subset M$). Karena $L$ membangkitkan ekstensi $LM|M$, bagian (ii) Definisi--Teorema \ref{def:pins} menunjukkan bahwa $LM|M$ murni tak separabel. Argumen pembangkit yang sama menunjukkan bahwa kompositum sebarang keluarga ekstensi murni tak separabel tetap murni tak separabel.
\end{proof}

\begin{lemma}\label{prop:pins-sep}
	Jika ekstensi aljabar $E|F$ separabel dan murni tak separabel, maka $[E:F]=1$.
\end{lemma}
\begin{proof}
	Karena ekstensi separabel dan ekstensi murni tak separabel sama-sama istimewa, cukup pertimbangkan kasus $E=F(x)$. Syaratnya setara dengan $[E:F]_s = [E:F]_i = 1$.
\end{proof}

Setiap ekstensi aljabar dapat diuraikan dalam dua tahap: mula-mula subekstensi separabel maksimal, lalu ekstensi murni tak separabel.
\begin{proposition}\label{prop:pins-devissage}
	Untuk setiap ekstensi aljabar $E|F$, misalkan subekstensi $E^\flat|F$ adalah penutup separabel $F$ di dalam $E$ (Definisi \ref{def:sep-closure}). Maka $E|E^\flat$ merupakan ekstensi murni tak separabel.
\end{proposition}
\begin{proof}
	Tuliskan polinomial minimal dari sembarang $x \in E$ sebagai $P_x \in F[X]$, lalu terapkan operasi \eqref{eqn:P-to-P-flat} untuk memperoleh $P_x(X) = P_x^\flat\left( X^{p^m} \right)$, dengan $P_x^\flat$ tak tereduksi dan tanpa akar ganda di penutup aljabar. Ini menunjukkan bahwa $y := x^{p^m}$ separabel atas $F$; karena $x^{p^m} \in F(y) \subset E^\flat$, unsur $x$ murni tak separabel atas $E^\flat$.
\end{proof}

Karena ekstensi murni tak separabel tertutup terhadap kompositum, kita juga dapat mendefinisikan subekstensi murni tak separabel maksimal $E^\natural$ dari $E|F$, yang disebut \emph{penutup murni tak separabel} $F$ di dalam $E$. Dibandingkan dengan hasil sebelumnya, wajar untuk bertanya apakah $E|E^\natural$ separabel. Berikut karakterisasi sederhananya.
\begin{proposition}\label{prop:pins-compositum}
	Dengan notasi di atas, ekstensi $E|E^\natural$ separabel $\iff E = E^\flat E^\natural$.
\end{proposition}
\begin{proof}
	($\implies$): Karena $E|E^\flat$ murni tak separabel, $E | E^\flat E^\natural$ juga demikian. Demikian pula, separabilitas $E|E^\natural$ menyiratkan separabilitas $E | E^\flat E^\natural$. Lemma \ref{prop:pins-sep} memberikan $E = E^\flat E^\natural$.
	
	($\impliedby$): Cukup perhatikan bahwa separabilitas $E^\flat|F$ menyiratkan separabilitas $E^\flat E^\natural | E^\natural$.
\end{proof}
Latihan akan memberikan contoh $E|E^\natural$ tak separabel. Setelah itu kita akan membuktikan bahwa ekstensi normal selalu memenuhi $E|E^\natural$ separabel (Teorema \ref{prop:normal-ext-structure}).

\section{Ekstensi Transenden}
Tetapkan medan $F$. Unsur dalam ekstensi $\Omega|F$ yang tidak aljabar disebut unsur transenden. Jadi $x \in \Omega$ transenden jika dan hanya jika homomorfisme $F$-aljabar
\begin{align*}
	F[X] & \longrightarrow \Omega \\
	X & \longmapsto x
\end{align*}
bersifat injektif; dalam hal ini $F[X] \rightiso F[x] \subsetneq F(x) \leftiso F(X)$. Sebagai contoh, $\pi, e \in \R$, yang dikenal dari analisis matematika, keduanya transenden atas $\Q$; lihat pembuktiannya dalam \cite[\S 7]{ZhP}. Terinspirasi oleh hal ini, kita memberikan definisi multivariabel berikut.

\begin{definition}\label{def:alg-indep}\index{daishuwuguan@bebas secara aljabar (algebraically independent)}
	Subhimpunan $\mathcal{X} \subset \Omega$ disebut \emph{bebas secara aljabar} atas $F$ jika memenuhi syarat berikut: untuk semua $n \geq 1$, unsur-unsur berbeda $x_1, \ldots, x_n \in \mathcal{X}$, dan polinomial $P \in F[X_1, \ldots, X_n]$, berlaku
	\[ P(x_1, \ldots, x_n)=0 \iff P=0. \]
\end{definition}
Dengan kata lain, unsur-unsur $\mathcal{X}$ tidak mempunyai relasi aljabar atas $F$ selain $0=0$; secara ekuivalen, $F[\mathcal{X}] \hookrightarrow \Omega$.

\begin{example}
	Menurut teorema dasar polinomial simetris, Teorema \ref{prop:fund-thm-symmetric-poly}, polinomial simetris elementer $e_0, \ldots, e_n$ bebas secara aljabar dalam medan fungsi rasional $F(X_1, \ldots, X_n)$. Submedan yang dibangkitkannya tidak lain adalah $F(X_1, \ldots, X_n)^{\mathfrak{S}_n}$.
\end{example}

Mulai sekarang, definisikan
\[ \text{Indep}_F(\Omega) := \left\{ \mathcal{X} \subset \Omega: \text{bebas secara aljabar} \}\right. \]
Menurut definisi, $\emptyset \in \text{Indep}_F(\Omega)$, dan $\mathcal{X} \in \text{Indep}_F(\Omega)$ menyiratkan bahwa setiap subhimpunan $\mathcal{X}$ juga termasuk dalam $\text{Indep}_F(\Omega)$. Himpunan $\text{Indep}_F(\Omega)$, dengan urutan $\subset$, merupakan poset tak kosong.

\begin{definition-theorem}\index{chaoyueji@basis transendensi (transcendence basis)}
	Setiap ekstensi $\Omega|F$ mempunyai subhimpunan maksimal yang bebas secara aljabar (yang boleh kosong); subhimpunan semacam itu bagi $\Omega|F$ disebut \emph{basis transendensi}.
\end{definition-theorem}
Jelas bahwa basis transendensi tak kosong ada jika dan hanya jika $\Omega|F$ bukan ekstensi aljabar.
\begin{proof}
	Menurut Teorema \ref{prop:Zorn}, cukup dibuktikan bahwa dalam poset $(\text{Indep}_F(\Omega), \subset)$ setiap rantai $\{ \mathcal{X}_i: i \in I\}$ mempunyai batas atas. Ambil $\mathcal{X} := \bigcup_{i \in I} \mathcal{X}_i$. Karena setiap $x_1, \ldots, x_n \in \mathcal{X}$ termuat dalam suatu $\mathcal{X}_i$, kebebasan aljabar tetap berlaku untuk $\mathcal{X}$.
\end{proof}

Misalkan $\mathcal{B}$ basis transendensi. Dari pembahasan di atas diketahui:
\begin{itemize}
	\item kebebasan aljabar setara dengan fakta bahwa subekstensi $F(\mathcal{B})|F$ yang dibangkitkan $\mathcal{B}$ secara alami isomorfik dengan medan fungsi rasional yang himpunan peubahnya $\mathcal{B}$;
	\item $\Omega|F(\mathcal{B})$ adalah ekstensi aljabar; jika tidak, suatu unsur transenden atas $F(\mathcal{B})$, yakni $x \in \Omega$, akan membuat $\mathcal{B} \cup \{x\} \in \text{Indep}_F(\Omega)$;
	\item sebaliknya, subhimpunan $\mathcal{B}$ yang memenuhi kedua sifat di atas pasti merupakan basis transendensi.
\end{itemize}

Suatu ekstensi medan dapat mempunyai beberapa basis transendensi, tetapi kardinalitasnya ditentukan secara unik. Sekarang kita membuktikan hal ini.
\begin{lemma}\label{prop:generator-cardinal-field}
	Misalkan $\mathcal{B}$ dan $\mathcal{B}'$ dua basis transendensi dari $\Omega|F$. Jika $\mathcal{B}$ tak hingga, maka $|\mathcal{B}'| \geq |\mathcal{B}|$.
\end{lemma}
\begin{proof}
	Setiap $b' \in \mathcal{B}'$ di atas $F(\mathcal{B})$ mempunyai polinomial minimal $P_{b'} \in F(\mathcal{B})[Y]$. Tulis koefisien-koefisiennya sebagai pecahan tereduksi dan definisikan himpunan berhingga $E_{b'} := \left\{ b \in \mathcal{B}: \text{muncul dalam koefisien $P_{b'}$} \right\}$. Ambil gabungan $\mathcal{B}'' := \bigcup_{b'} E_{b'} \subset \mathcal{B}$. Maka $\Omega|F(\mathcal{B}'')$ aljabar, sehingga harus berlaku $\mathcal{B}''=\mathcal{B}$. Jadi $\mathcal{B}'$ juga tak hingga. Korolari \ref{prop:cardinal-max} memberikan
	\[ \max\{|\mathcal{B}'|, \aleph_0 \} = |\mathcal{B}'| \cdot \aleph_0 \geq \left| \bigcup_{b' \in \mathcal{B}'} E_{b'} \right| = |\mathcal{B}|, \]
	Paling kiri tidak lain adalah $|\mathcal{B}'|$.
\end{proof}

\begin{lemma}[Sifat pertukaran]\label{prop:exchange-field}
	Misalkan $\mathcal{B}$ dan $\mathcal{B}'$ dua basis transendensi berhingga dari $\Omega|F$, dan $b' \in \mathcal{B}' \smallsetminus \mathcal{B}$. Maka terdapat $b \in \mathcal{B} \smallsetminus \mathcal{B}'$ sehingga $(\mathcal{B}' \smallsetminus \{b'\}) \cup \{b\}$ tetap merupakan basis transendensi.
\end{lemma}
\begin{proof}
	Karena $\mathcal{B}'$ tak kosong, $\mathcal{B}$ juga tak kosong. Setiap $b \in \mathcal{B}$ aljabar atas $F(\mathcal{B}')$, sehingga merupakan akar suatu polinomial $P_b \in F[\mathcal{B}'][Y]$. Kita nyatakan bahwa $b'$ harus muncul dalam suatu $P_b$. Jika tidak, setiap $b \in \mathcal{B}$ aljabar atas $F(\mathcal{B}' \smallsetminus \{b'\})$, sehingga $\mathcal{B}' \smallsetminus \{b'\}$ merupakan basis transendensi, sebuah kontradiksi.
	
	Ambil $b \in \mathcal{B}$ sedemikian sehingga $b'$ muncul dalam $P_b$. Definisikan $\mathcal{B}'' := (\mathcal{B}' \smallsetminus \{b'\}) \cup \{b\}$.
	\begin{compactenum}[(a)]
		\item Setiap $\beta' \in \mathcal{B}'$ aljabar atas $F(\mathcal{B}'')$: cukup gunakan $P_b(b)=0$ untuk memeriksa kasus $\beta' = b'$.
		\item Jika $\mathcal{B}'' \notin \text{Indep}_F(\Omega)$, karena $\mathcal{B}' \smallsetminus \{b'\} \in \text{Indep}_F(\Omega)$, unsur $b$ aljabar atas $F(\mathcal{B}' \smallsetminus \{b'\})$. Dari pengamatan sebelumnya, $b'$ juga aljabar atas $F(\mathcal{B}' \smallsetminus \{b'\})$, sebuah kontradiksi.
	\end{compactenum}
	Jadi $\mathcal{B}''$ memang basis transendensi. Terakhir, perhatikan bahwa $b \notin \mathcal{B}'$. Jika tidak, maka $b' \notin \mathcal{B} \implies b' \neq b \implies \mathcal{B}'' = \mathcal{B}' \smallsetminus \{b'\}$, bertentangan dengan definisi basis transendensi sebagai subhimpunan bebas secara aljabar yang maksimal.
\end{proof}

\begin{definition-theorem}\index{yukuozhang! derajat transendensi (transcendence degree)}\index[sym1]{trdeg@$\trdeg$}
	Dalam ekstensi medan $\Omega|F$, setiap dua basis transendensi $\mathcal{B}$, $\mathcal{B}'$ memenuhi $|\mathcal{B}| = |\mathcal{B}'|$. Kardinal bersama ini bagi $\Omega|F$ disebut \emph{derajat transendensi} dan ditulis $\trdeg(\Omega|F)$.
\end{definition-theorem}
\begin{proof}
	Gagasannya sama dengan Definisi--Teorema \ref{def:dimension-vector-space}. Jika terdapat basis transendensi tak hingga, Lemma \ref{prop:generator-cardinal-field} menyiratkan bahwa semua basis transendensi juga tak hingga dan $|\mathcal{B}| \leq |\mathcal{B}'| \leq |\mathcal{B}|$, sehingga $|\mathcal{B}|=|\mathcal{B}'|$.
		
	Andaikan semua basis transendensi berhingga. Lemma \ref{prop:exchange-field} menyiratkan bahwa seluruh basis transendensi dalam $\Omega$ membentuk matroid seperti dalam Definisi \ref{def:matroid}; Proposisi \ref{prop:matroid-rank} langsung memberikan $|\mathcal{B}|=|\mathcal{B}'|$.\index{nizhen}
\end{proof}

Jika $\Omega$ tertutup secara aljabar, pembahasan di atas menunjukkan bahwa $\Omega$ adalah penutup aljabar dari $F(\mathcal{B})$, sedangkan $F(\mathcal{B})$ adalah medan fungsi rasional dengan himpunan peubah $\mathcal{B}$. Korolari berikut menjadi jelas.
\begin{corollary}
	Misalkan $\Omega_1, \Omega_2$ ekstensi dari $F$, dan $\Omega_1, \Omega_2$ keduanya tertutup secara aljabar. Maka terdapat isomorfisme $F$-aljabar $\Omega_1 \simeq \Omega_2$ jika dan hanya jika $\trdeg(\Omega_1|F) = \trdeg(\Omega_2|F)$.
\end{corollary}

\begin{corollary}\label{prop:ACF-categoricity}
	Misalkan $E_1$, $E_2$ adalah medan-medan tertutup secara aljabar, $\mathrm{char}(E_1)=\mathrm{char}(E_2)$ dan $\kappa := |E_1| = |E_2| > \aleph_0$. Maka terdapat isomorfisme medan $E_1 \simeq E_2$.
\end{corollary}
\begin{proof}
	Misalkan $F$ submedan prima bersama dari $E_1$ dan $E_2$, sehingga $|F| \leq \aleph_0$. Ambil basis transendensi $\mathcal{B}_i \subset E_i$ (yang harus tak hingga), dan perhatikan bahwa Proposisi \ref{prop:alg-ext-cardinality} memberikan $\kappa = |F(\mathcal{B}_i)| = |F[\mathcal{B}_i]|$. Korolari \ref{prop:cardinal-max} memberikan
	\begin{align*}
		|\mathcal{B}_i| & \leq |F[\mathcal{B}_i]| \leq \sum_{n \geq 0} \left| \{P \in F[\mathcal{B}_i]: \deg P = n \} \right| \\
		& \leq \sum_{n \geq 0} |\mathcal{B}_i| = |\mathcal{B}_i|.
	\end{align*}
	Dengan demikian $\kappa = |\mathcal{B}_i|$, kemudian menggunakan kesimpulan sebelumnya didapatlah $E_1 \simeq E_2$.
\end{proof}
Sebagai kasus khusus, semua medan tertutup secara aljabar berkarakteristik nol dan berkardinalitas $2^{\aleph_0} = |\CC|$ (Contoh \ref{eg:continuum}) isomorfik dengan $\CC$; ini merupakan trik kecil yang sering digunakan dalam geometri aljabar aritmetika. Dalam istilah teori model, korolari ini mengatakan bahwa teori medan tertutup secara aljabar berkarakteristik $p \geq 0$, yaitu $\text{ACF}_p$, bersifat $\kappa$-kategoris, dengan $\kappa$ sembarang kardinal $> \aleph_0$; lihat \cite[\S 6.1.1 dan Bab 8]{Feng17}. \index{kappa-dengshitonggou@$\kappa$-kategorisitas ($\kappa$-categoricity)} \index{moxinglun@teori model}

\begin{proposition}
	Jika $\mathcal{C}$ basis transendensi dari $L|E$ dan $\mathcal{B}$ basis transendensi dari $E|F$, maka $\mathcal{B} \sqcup \mathcal{C}$ adalah basis transendensi dari $L|F$. Khususnya, $\trdeg(L|F) = \trdeg(L|E) + \trdeg(E|F)$.
\end{proposition}
\begin{proof}
	Dari definisi, kebebasan aljabar $\mathcal{C}$ atas medan fungsi rasional $F(\mathcal{B}) \subset E$ menyiratkan kebebasan $\mathcal{B} \sqcup \mathcal{C}$ atas $F$. Keistimewaan ekstensi aljabar (Proposisi \ref{prop:alg-ext-dist}) menyiratkan bahwa $E(\mathcal{C}) | F(\mathcal{B})(\mathcal{C}) = F(\mathcal{B} \sqcup \mathcal{C})$ aljabar, sehingga $L|F(\mathcal{B} \sqcup \mathcal{C})$ juga aljabar.
\end{proof}

\section{Aplikasi Produk Tensor}\label{sec:tensor-field}
Pertimbangkan ekstensi medan $\Omega|F$ dan subekstensi $E|F$ serta $E'|F$. Sifat universal dari produk tensor memberikan homomorfisme dari $F$-aljabar
\begin{align*}
	m: E \dotimes{F} E' & \longrightarrow EE' \subset \Omega \\
	x \otimes y & \longmapsto xy.
\end{align*}
Fakta sederhana: $E \dotimes{F} E' \neq \{0\}$ (catatan \ref{rem:algebra-otimes-zero}).

\begin{lemma}
	Kernel dari homomorfisme $m$ adalah ideal prima, $EE'$ dapat disamakan dengan medan pecahan dari $\Image(m)$; ketika $E|F$ dan $E'|F$ adalah ekstensi aljabar, $\Image(m)=EE'$.
\end{lemma}
\begin{proof}
	Pernyataan pertama mengikuti dari fakta bahwa $\Image(m) \subset \Omega$ merupakan gelanggang integral tak nol. Deskripsi eksplisit kompositum medan menunjukkan bahwa $EE'$ tepat merupakan medan pecahan dari $\Image(m)$. Jika $E|F$ dan $E'|F$ ekstensi aljabar, maka $EE'|F$ juga aljabar. Karena itu, invers setiap unsur tak nol $x \in \Image(m)$ berada dalam $F[x]$, sehingga $\Image(m)$ sendiri sudah merupakan medan.
\end{proof}

\begin{definition}\index{xianxingwujiao@saling bebas linear (linearly disjoint)}
	Jika homomorfisme $m$ di atas injektif, maka $E|F$ dan $E'|F$ disebut \emph{saling bebas linear}.
\end{definition}
Sifat saling bebas linear mempunyai karakterisasi elementer dalam aljabar linear berikut; ini sekaligus menunjukkan bahwa, seperti banyak sifat lain dalam teori medan, sifat tersebut pada hakikatnya ``berhingga''.

\begin{proposition}
	Sifat-sifat berikut setara:
	\begin{compactenum}[(i)]
		\item subekstensi $E|F$ dan $E'|F$ saling bebas linear;
		\item setiap keluarga berhingga unsur $x_1, \ldots, x_n \in E$ yang bebas linear atas $F$ tetap bebas linear atas $E'$;
		\item setiap keluarga berhingga unsur $y_1, \ldots, y_n \in E'$ yang bebas linear atas $F$ tetap bebas linear atas $E$.
	\end{compactenum}
\end{proposition}
\begin{proof}
	Berdasarkan simetri, cukup buktikan (i) $\iff$ (ii). Ambil unsur-unsur bebas linear $x_1, \ldots, x_n \in E$, dan pilih subruang-$F$ $W \subset E$ sedemikian sehingga $\lrangle{x_1, \ldots, x_n} \oplus W = E$. Hasil kali tensor mempertahankan jumlah langsung, sehingga panah $i$ dalam diagram berikut injektif:
	\[\begin{tikzcd}
		\lrangle{x_1, \ldots, x_n} \dotimes{F} E' \arrow[hookrightarrow, r, "i", "\text{suku langsung}"' inner sep=0.6em] & \left( \lrangle{x_1, \ldots, x_n} \oplus W \right) \dotimes{F} E' \arrow[equal, r] & E \dotimes{F} E' \arrow[d, "m"] \\
		{E'}^{\oplus n} \arrow[u, "\simeq"] \arrow[rr, "{(a_1, \ldots, a_n) \mapsto \sum_{i=1}^n a_i x_i}"'] & & EE'
	\end{tikzcd}\]
	Mudah diverifikasi bahwa diagram itu komutatif. Jika (i) berlaku, $m$ injektif, sehingga panah horizontal bawah juga injektif; ini tepat merupakan (ii). Sebaliknya, andaikan (ii). Menurut konstruksi hasil kali tensor, setiap unsur $E \dotimes{F} E'$ berada dalam suatu subruang vektor berbentuk $\lrangle{x_1, \ldots, x_n} \dotimes{F} E'$, dan kita boleh menganggap $x_1, \ldots, x_n$ bebas linear. Dari diagram di atas, pembatasan $m$ pada $\lrangle{x_1, \ldots, x_n} \dotimes{F} E'$ injektif. Jadi $m$ injektif.
\end{proof}

Aplikasi sederhana lain dari hasil kali tensor: setiap dua ekstensi medan dapat dibentuk kompositumnya di dalam suatu ekstensi medan yang lebih besar.
\begin{proposition}
	Untuk setiap medan $F$ dan ekstensi medan $F_1|F$, $F_2|F$, selalu terdapat ekstensi medan $\Omega|F$ dan $F$-pembenaman $\iota_i: F_i \hookrightarrow  \Omega$ ($i=1,2$). Khususnya, kompositum $F_1 F_2$ di dalam $\Omega$ bermakna.
\end{proposition}
\begin{proof}
	Telah disebutkan sebelumnya bahwa $F_1 \dotimes{F} F_2$ merupakan $F$-aljabar tak nol. Proposisi \ref{prop:existence-maximal-ideal} menyatakan bahwa aljabar ini mempunyai ideal maksimal $\mathfrak{m}$. Maka kita dapat mengambil $\Omega := (F_1 \dotimes{F} F_2) \big/ \mathfrak{m}$ dan homomorfisme $F$-aljabar $\iota_i: F_i \hookrightarrow F_1 \dotimes{F} F_2 \twoheadrightarrow \Omega$ ($i=1,2$).
\end{proof}

Berikutnya tetapkan medan $F$ dan penutup aljabarnya $\overline{F}$.
\begin{definition}\label{def:etale-alg}\index{daishu!étale (étale)}\index{daishu!dapat didiagonalisasi (diagonalizable)}
	Misalkan $A$ adalah $F$-aljabar komutatif yang tidak nol.
	\begin{itemize}
		\item Jika terdapat isomorfisme $F$-aljabar $A \simeq \underbracket{F \times \cdots \times F}_{n\;\text{suku}} = F^n$, dengan $n = \dim_F A$, maka $A$ disebut dapat didiagonalisasi;
		\item Jika $\overline{F}$-aljabar $A \dotimes{F} \overline{F}$ (perubahan basis, lihat \S\ref{sec:algebra-tensor-product}) dapat didiagonalisasi, maka $A$ disebut \emph{étale}.
	\end{itemize}
\end{definition}

Istilah étale berasal dari geometri aljabar; salah satu kelebihannya ialah ketertutupan terhadap perubahan basis. Sebelum membuktikan pernyataan berikut, perhatikan dua fakta sederhana: hasil kali tensor aljabar-aljabar yang dapat didiagonalisasi juga dapat didiagonalisasi, demikian pula perubahan basisnya.
\begin{lemma}\label{prop:etale-algebra-basechange}
	Jika $A, B$ adalah $F$-aljabar étale, maka $A \dotimes{F} B$ juga étale. Untuk setiap ekstensi medan $L|F$ dan $F$-aljabar étale $A$, $L$-aljabar $A \dotimes{F} L$ juga étale.
\end{lemma}
\begin{proof}
	Untuk pernyataan pertama, proposisi \ref{prop:algebra-base-change-monoidal} menyediakan isomorfisme alami dari $\overline{F}$-aljabar
	\[ (A \dotimes{F} B) \dotimes{F} \overline{F} \simeq (A \dotimes{F} \overline{F}) \dotimes{\overline{F}} (B \dotimes{F} \overline{F}), \]
	Dengan demikian, masalah pertama direduksi ke kasus $F=\overline{F}$ dengan $A$ dan $B$ dapat didiagonalisasi. Untuk pernyataan kedua, bagi $L$, ambil penutup aljabarnya $\overline{L}$; maka $\overline{F}$ dapat dibenamkan ke dalam $\overline{L}$. Isomorfisme alami dalam Proposisi \ref{prop:tensor-unit} memberikan
	\[ (A \dotimes{F} L) \dotimes{L} \overline{L} \simeq A \dotimes{F} \overline{L} \simeq (A \dotimes{F} \overline{F}) \dotimes{\overline{F}} \overline{L} \]
	Ini mereduksi masalah ke kasus $F=\overline{F}$ dengan $A$ dapat didiagonalisasi, yang tercakup dalam pengamatan sebelumnya.
\end{proof}

Untuk setiap ekstensi medan $E|F$, memberikan isomorfisme $E$-aljabar $\Lambda: A_E := A \dotimes{F} E \rightiso E^n$ setara dengan memberikan keluarga $\lambda_1, \ldots, \lambda_n \in \Hom_{E\dcate{Alg}}(A_E, E)$ yang menjadi basis ruang vektor dual $A_E^\vee := \Hom_E(A_E, E)$; korespondensinya ditentukan oleh $\Lambda = (\lambda_1, \ldots, \lambda_n)$. Karena $\overline{F}|F$ adalah gabungan subekstensi berhingga, isomorfisme $A_{\overline{F}} \rightiso \overline{F}^n$ bagi aljabar étale $A$ selalu dapat didefinisikan pada suatu subekstensi berhingga $E|F$ yang cukup besar; jadi $A_E$ dapat didiagonalisasi.

\begin{lemma}\label{prop:etale-algebra-subalgebras}
	$F$-aljabar étale $A$ hanya mempunyai berhingga banyak subaljabar dan ideal; subaljabar serta hasil bagi oleh ideal sejati tetap étale. Lebih tepat, jika $A \simeq F^n$ dapat didiagonalisasi:
	\begin{itemize}
		\item Subaljabar berkorespondensi satu-satu dengan dekomposisi disjungsi $\{1, \ldots, n\} = I_1 \sqcup \cdots \sqcup I_r$, subaljabar yang sesuai ditentukan oleh unsur idempoten
			\[ e_I := (x_i)_{i=1}^n, \quad x_i = \begin{cases} 1, & i \in I \\ 0, & i \notin I \end{cases}, \quad I = I_1, \ldots, I_r \]
			membangkitkannya, sehingga subaljabar tersebut juga dapat didiagonalisasi;
		\item ideal-ideal berkorespondensi satu-satu dengan subhimpunan $I \subset \{1, \ldots, n\}$; ideal yang bersesuaian dibangkitkan oleh $e_I$, dan jika $I \neq \{1, \ldots, n\}$, aljabar hasil baginya juga dapat didiagonalisasi.
	\end{itemize}
\end{lemma}
\begin{proof}
	Pertama tangani kasus $A = F^n$. Pernyataan tentang ideal dan hasil bagi relatif mudah; kita fokus pada subaljabar $B \subset A$. Pembatasan proyeksi koordinat $e_i: A \to F$ ($i=1, \ldots, n$) pada $B$ memberikan pembangkit-pembangkit bagi $B^\vee = \Hom_F(B,F)$, yaitu $\bar{e}_i \in \Hom_{F\dcate{Alg}}(B, F)$. Pemilihan basis $\bar{e}_{i_1}, \bar{e}_{i_2}, \ldots, \bar{e}_{i_r}$ memberikan $B \rightiso F^r$, sehingga $B$ dapat didiagonalisasi. Sekarang definisikan $f_j$ sebagai prapeta $(0, \ldots, \underbracket{1}_{\text{suku ke-$j$}}, \ldots, 0) \in F^r$ di bawah $B \rightiso F^r$. Maka $B = F[f_1, \ldots, f_r]$ dan
	\begin{gather*}
		f_j^2 = f_j, \quad \sum_{j=1}^r f_j = 1, \\
		j \neq k \implies f_j f_k = 0;
	\end{gather*}
	(bandingkan kasus modul dalam \S\ref{sec:indecomposable-mod}). Dengan memeriksa persamaan ini per koordinat dalam $A = F^n$, terlihat bahwa $f_1, \ldots, f_r$ harus berasal dari dekomposisi $\{1, \ldots, n\} = I_1 \sqcup \cdots \sqcup I_r$ dalam pernyataan.
	
	Dalam kasus étale, $- \dotimes{F} \overline{F}$ memetakan subaljabar (atau ideal) $A$ ke subaljabar (atau ideal) $A \dotimes{F} \overline{F} \simeq \overline{F}^n$, dan peta ini injektif: untuk setiap subruang-$F$ $B \subset A$, berlaku $(B \dotimes{F} \overline{F}) \cap (A \otimes 1) = B \otimes 1 = B$. Pernyataan selebihnya direduksi ke kasus yang dapat didiagonalisasi.
\end{proof}

Kita memperoleh karakterisasi ringkas ekstensi berhingga separabel.
\begin{proposition}\label{prop:etale-algebra-separable}
	Ekstensi berhingga $E|F$ separabel jika dan hanya jika $E$ merupakan $F$-aljabar étale. Dalam hal ini terdapat isomorfisme $F^\mathrm{sep}$-aljabar $E \dotimes{F} F^\mathrm{sep} \simeq (F^\mathrm{sep})^{[E:F]}$.
\end{proposition}
\begin{proof}
	Andaikan $E$ étale. Proposisi \ref{prop:pins-devissage} menguraikan $E|F$ menjadi ekstensi murni tak separabel $E|E^\flat$ dan ekstensi separabel $E^\flat|F$. Untuk membuktikan bahwa $E|F$ separabel, kita boleh menganggap $p := \text{char}(F) > 0$. Terapkan Lemma \ref{prop:etale-algebra-subalgebras}: di bawah isomorfisme $E \dotimes{F} \overline{F} \rightiso \overline{F}^n$, $E^\flat \dotimes{F} \overline{F}$ bersesuaian dengan subaljabar $B$ yang ditentukan oleh gabungan lepas $\{1, \ldots, n\} = I_1 \sqcup \cdots \sqcup I_r$. Jika $x \in E \dotimes{F} \overline{F}$ bersesuaian dengan $y \in \overline{F}^n$, maka terdapat $m \gg 0$ dengan $y^{p^m} \in B$. Karena setiap unsur $\overline{F}$ mempunyai tepat satu akar pangkat $p^m$, deskripsi eksplisit $B$ memberikan $y \in B$, sehingga $x \in E^\flat \dotimes{F} \overline{F}$. Karena kita bekerja di atas medan, $E^\flat \dotimes{F} \overline{F} = E \dotimes{F} \overline{F}$ menyiratkan $E^\flat = E$. Jadi $E|F$ separabel.
	
	Sekarang andaikan $E|F$ separabel dan tetapkan $n := [E:F]$. Menurut teorema unsur primitif, Teorema \ref{prop:prim-element-separable}, kita boleh menganggap $E = F(u) \simeq F[X]/(P_u)$, dengan $P_u$ polinomial minimal $u$. Karena di atas $\overline{F}$ berlaku $P_u = \prod_{i=1}^n (X-\alpha_i)$ tanpa faktor berulang, tetapkan $L := F(\alpha_1, \ldots, \alpha_n) \subset F^\text{sep}$. Teorema Sisa Cina \ref{prop:CRT} memberikan isomorfisme $L$-aljabar
	\[ E_L := E \dotimes{F} L \simeq L[X] \bigg/ ( \prod_{i=1}^n (X-\alpha_i) ) = \prod_{i=1}^n L[X]/(X-\alpha_i) \simeq L^n, \]
	Sehingga diketahui $E \dotimes{F} L$ dapat didiagonalisasi; maka $F^\mathrm{sep}$-aljabar $E \dotimes{F} F^\mathrm{sep}$ juga dapat didiagonalisasi, jadi $E$ étale. Bagian kedua dari pernyataan terbukti.
\end{proof}
\begin{remark}
	Dengan demikian, bagian pertama Teorema \ref{prop:prim-element-separable} dapat dipandang sebagai kasus khusus Teorema \ref{prop:prim-element}: ekstensi separabel $E|F$ étale, sehingga Lemma \ref{prop:etale-algebra-subalgebras} menyiratkan bahwa $E|F$ hanya mempunyai berhingga banyak medan antara.
\end{remark}

Hubungan separabilitas dengan hasil kali tensor dapat diperluas ke ekstensi tak hingga, bahkan ekstensi nonaljabar. Pekerjaan utama di bidang ini dilakukan oleh Mac Lane; kuncinya ialah menggunakan basis-$p$ yang diperkenalkan Teichmüller. Teori ini cukup rumit dan termasuk cabang teori gelanggang komutatif, sehingga lebih baik dibahas pada kesempatan lain.

Jika suatu gelanggang tidak mempunyai unsur nilpoten tak nol, gelanggang itu disebut \emph{tereduksi}. Jelas bahwa subgelanggang dari gelanggang tereduksi juga tereduksi.
\begin{corollary}\label{prop:etale-alg-characterization}
	Untuk $F$-aljabar komutatif tak nol berdimensi berhingga $A$, kondisi berikut setara.
	\begin{compactenum}[(i)]
		\item $A$ adalah aljabar étale.
		\item untuk setiap ekstensi medan $E|F$, aljabar $A_E$ tereduksi;
		\item aljabar $A_{\overline{F}}$ tereduksi;
		\item terdapat ekstensi berhingga separabel $E_1, \ldots, E_n$ dari $F$ ($n \geq 1$) sedemikian sehingga $A \simeq E_1 \times \cdots \times E_n$.
	\end{compactenum}
\end{corollary}
\begin{proof}
	(i) $\implies$ (ii): Dari $A \hookrightarrow A \dotimes{F} \overline{F} \simeq \overline{F}^{\dim A}$ diketahui bahwa $A$ tereduksi. Karena $A_E$ merupakan $E$-aljabar étale, argumen yang sama menunjukkan bahwa $A_E$ tereduksi.

	(ii) $\implies$ (iii) jelas; sekarang buktikan (iii) $\implies$ (iv). Jika $A$ medan, maka $A|F$ separabel: untuk setiap $u \in A$, submedan $F[u] \simeq F[X]/(P_u)$ menjadi tereduksi sesudah perubahan basis ke $\overline{F}$, sehingga $P_u$ tidak mempunyai akar ganda dalam $\overline{F}$. Jika $A$ bukan medan, terdapat ideal sejati tak nol $\mathfrak{a} \subsetneq A$. Karena $A$ tereduksi, $\{0\} \neq \mathfrak{a}^2 \subset \mathfrak{a}$; pilih $\mathfrak{a}$ berdimensi minimal, sehingga $\mathfrak{a}^2 = \mathfrak{a}$. Akan dibuktikan bahwa terdapat idempoten $e \in A$ dengan $\mathfrak{a} = Ae$. Maka ada dekomposisi produk langsung nontrivial $A = eA \times (1-e)A$; rekursi pada $\dim_F A$ menguraikan $A$ sebagai produk langsung ekstensi-ekstensi separabel.

	Ambil $\mathfrak{a}$ sebagai pembangkit $a_1, \ldots, a_m$, maka $\mathfrak{a}^2 = \mathfrak{a}$ menyiratkan keberadaan matriks $B \in M_m(\mathfrak{a})$ sehingga $(a_1, \ldots, a_m) B = (a_1, \ldots, a_m)$, yaitu $(a_1, \ldots, a_m) \left( 1_{m \times m} - B \right) = 0$. Dengan menggunakan hubungan antara matriks adjoint dan determinan (teorema \ref{prop:matrix-det}), selanjutnya diturunkan
	\[ \det(1-B) (a_1, \ldots, a_m) = (a_1, \ldots, a_m) \left( 1_{m \times m} - B \right) \left( 1_{m \times m} - B \right)^\vee = 0. \]
	Jadi $\forall i,\; \det(1-B)a_i = 0$ (teknik serupa pernah digunakan ketika membuktikan Teorema \ref{prop:integrality-finiteness}). Namun $\det(1-B)$ juga dapat dinyatakan dalam bentuk $1-e$, dengan $e \in \mathfrak{a}$; mudah diverifikasi bahwa $\forall a \in \mathfrak{a},\; ea = a$, sehingga $Ae = \mathfrak{a}$ dan $e$ adalah unsur idempoten tak nol. Bukti selesai.
	
	(iv) $\implies$ (i). Produk langsung aljabar-aljabar étale tetap étale. Proposisi \ref{prop:etale-algebra-separable} langsung menyatakan bahwa jika $E_1|F, \ldots, E_n|F$ berhingga separabel, maka $E_1 \times \cdots \times E_n$ étale.
\end{proof}

Kini struktur setiap aljabar étale sudah sepenuhnya jelas. Pembaca mungkin bertanya: jika aljabar étale hanya produk ekstensi separabel, mengapa konsep ini diperkenalkan? Salah satu alasannya ialah perubahan basis mempertahankan sifat étale (Lemma \ref{prop:etale-algebra-basechange}); pencarian sifat ketertutupan semacam ini merupakan gagasan dasar matematika modern. Alasan lain ialah kategori $F$-aljabar étale ekuivalen dengan kategori himpunan-$\Gamma_F$ berhingga, dengan $\Gamma_F$ grup Galois absolut $F$ yang dilengkapi topologi Krull. Konsep terkait akan dijelaskan secara rinci dalam Bab Sembilan dan latihannya.

\begin{Exercises}
	\item Buktikan bahwa untuk setiap ekstensi $\Omega|F$ dan unsur aljabar $\alpha \in \Omega$, jika $[F(\alpha):F]$ ganjil, maka $F(\alpha)=F(\alpha^2)$.
	\item Misalkan $\alpha, \beta \in \overline{F}$ adalah akar dari polinomial monik tak tereduksi $P \in F[X]$. Buktikan bahwa ketika $\deg P > 1$ adalah bilangan ganjil, $\alpha+\beta \notin F$.
	\begin{hint} Jika $c := \alpha + \beta \in F$, maka $P(X)$ dan $Q(X) := (-1)^{\deg P} P(c-X)$ memiliki akar bersama, dari sini membuktikan $P=Q$, kemudian memasangkan akar sesuai $x \leftrightarrow c-x$.\end{hint}
	\item Untuk setiap medan $F$ berkarakteristik $p>0$, buktikan bahwa penutup sempurna $K := F^{p^{-\infty}}$ memenuhi $K^p = K$ (yakni $K$ medan sempurna). Jika $E$ medan sempurna dan $E|F$ ekstensi aljabar, buktikan bahwa $K|F$ dapat dibenamkan sebagai subekstensi $E|F$.
	\item (G.\ Elencwajg) tetapkan $\F_2 := \Z/2\Z$, ambil medan fungsi rasional dua variabel $F := \F_2(u,v)$ dan polinomial $P := X^6 + uvX^2 + u \in F[X]$. Buktikan:
		\begin{compactenum}[(i)]
			\item $P$ tak tereduksi dan $E := F[X]/(P)$ memenuhi $[E:F]=6$, $[E:F]_s = 3$; \begin{hint} Gunakan kriteria Eisenstein (Teorema \ref{prop:Eisenstein-criterion}) untuk membuktikan $P$ tak tereduksi, lalu hitung derajat separabel dengan Lemma \ref{prop:field-sdegree-F(x)}. \end{hint}
			\item penutup murni tak separabel $E^\natural \subset E$ sama dengan $F$. \begin{hint} Cukup buktikan $\beta \in E, \; \beta^2 \in F \implies \beta \in F$. Kembangkan semua suku dalam basis $1, X, \ldots, X^5 \; \bmod P$ dan periksa kapan $\beta^2 \in F$. \end{hint}
		\end{compactenum}
		Jadi dalam contoh ini $E|E^\natural$ tak separabel.
	\item Misalkan $p$ bilangan prima dan $\F_p := \Z/p\Z$. Perhatikan medan fungsi rasional dua peubah $K := \F_p(X,Y)$.
		\begin{enumerate}[(i)]
			\item Verifikasikan $K^p = \F_p(X^p, Y^p)$, $K=K^p(X,Y)$, dan $a \in K \implies [K^p(a) : K^p] \leq p$.
			\item Simpulkan dari Proposisi \ref{prop:field-tower-degree} dan \ref{prop:field-tower-sdegree} bahwa
				\begin{gather*}
					[K:K^p] = [\F_p(X,Y):\F_p(X^p,Y)] \cdot [\F_p(X^p,Y):\F_p(X^p,Y^p)] = p^2, \\
					[K:K^p]_s = [\F_p(X,Y):\F_p(X^p,Y)]_s \cdot [\F_p(X^p,Y):\F_p(X^p,Y^p)]_s = 1.
				\end{gather*}
			\item Verifikasi bahwa $\{ X^i Y^j: 0 \leq i,j < p \}$ membentuk basis ruang vektor-$K^p$ dari $K$.
			\item Membuktikan bahwa saat $c$ menelusuri unsur-unsur $K^p$ (tak terhingga banyaknya), setiap medan menengah $K^p(cX + Y)$ berbeda satu sama lain.
			\begin{hint} Jika $c \neq c'$ memberikan medan antara yang sama $M$, maka $(c-c')X \in M$, sehingga $X,Y \in M$ dan dengan demikian $M = K$, sebuah kontradiksi. \end{hint}
		\end{enumerate}
		Jadi $K|K^p$ adalah ekstensi berhingga murni tak separabel dengan tak hingga banyak medan antara, dan bukan ekstensi sederhana; bandingkan dengan Teorema \ref{prop:prim-element} dan \ref{prop:prim-element-separable}.
	\item Untuk medan fungsi rasional satu peubah $F(X)$ atas sembarang medan $F$, ambil $u := \frac{r}{s} \in F(X) \smallsetminus F$, dengan $r,s \in F[X]$ relatif prima.
		\begin{compactenum}[(i)]
			\item membuktikan bahwa $X$ adalah unsur aljabar di atas $F(u)$, $u$ adalah unsur transenden di atas $F$;
			\item Perkenalkan peubah bebas $Z, W$. Buktikan bahwa $P(Z,W) := r(Z) - Ws(Z)$ tak tereduksi sebagai unsur $F[Z,W]$, lalu bahwa $P(Z,u)$ tak tereduksi sebagai unsur $F(u)[Z]$;
			\item simpulkan $[F(X):F(u)] = \max\{\deg r, \deg s\}$.
		\end{compactenum}
	\item Berdasarkan yang di atas, membuktikan $\Aut_F(F(X))^{\mathrm{op}} \simeq \PGL(2,F) := \GL(2,F)/F^\times$, cara kerjanya adalah $X \cdot \bigl( \begin{smallmatrix} a & b \\ c & d \end{smallmatrix} \bigr) = \frac{aX+b}{cX+d}$.
	\item (Teorema Lüroth) Buktikan bahwa subekstensi $K|F$ dari $F(X)|F$ adalah $F$ atau berbentuk $K=F(u)$ seperti dalam soal sebelumnya. \index{Lüroth dingli@Teorema Lüroth}
	\begin{hint}
		Misalkan di atas $K \neq F$, peubah $X$ memiliki polinomial minimal $Q(Z) = Z^n + u_{n-1} Z^{n-1} + \cdots + u_0$; ambil $i$ sehingga $u = \frac{r}{s} = u_i \notin F$. Kita memiliki $Q(Z) \mid P(Z,u)$. Buktikan $P(Z,u) \in F^\times Q(Z)$ untuk mendapatkan $K=F(u)$.
	\end{hint}
	\item Tentukan $\Nm_{E|F}(E^\times) \subset F^\times$ untuk ekstensi berhingga $E|F$ berikut:
		\begin{compactenum}[(i)]
			\item $E$, $F$ adalah medan berhingga; \hint{Kita boleh menggunakan teori Galois.}
			\item $E=\CC$, $F=\R$;
			\item $E=\Q(\sqrt{-1})$, $F=\Q$.
		\end{compactenum}
	\item Menyambung dari sebelumnya, untuk medan berhingga $E, F$ buktikan $\Tr_{E|F}(E)=F$.
	\item Misalkan $F$ adalah medan berhingga, $a, b \in F^\times$, $c \in F$, buktikan bahwa ada $(x,y) \in F^2$ sedemikian sehingga $ax^2 + by^2 = c$.
	\item Misalkan $\sigma \in \Aut(\R)$.
		\begin{enumerate}[(i)]
			\item membuktikan $x \geq 0 \iff \sigma(x) \geq 0$. \hint{$x \geq 0 \iff \exists y \in \R, \; x = y^2$}
			\item Dari sini dapat disimpulkan bahwa untuk semua bilangan rasional $a < b$ berlaku $\sigma([a,b]) = [a,b]$.
			\item membuktikan $\sigma = \identity_{\R}$.
		\end{enumerate}
	\item membuktikan bahwa $\Aut(\CC)$ tak terhingga.
	\item Buktikan bahwa setiap ekstensi yang dibangkitkan secara berhingga dari $\Q$ dapat dibenamkan ke dalam $\CC|\Q$.
	\item Bangun medan tertutup secara aljabar $F_1, F_2$ berkarakteristik $0$ yang tidak isomorfik satu sama lain dan memenuhi $|F_1|=|F_2|=\aleph_0$.
	\item Misalkan $A$ aljabar komutatif berdimensi berhingga atas medan $F$. Buktikan bahwa $A$ étale jika dan hanya jika diskriminan dalam Definisi \ref{def:discriminant}, yaitu $d_A \neq 0$.
	\begin{hint}
		Diskriminan tidak nol jika dan hanya jika bentuk teras $(x,y) \mapsto \Tr_{A|F}(xy)$ tak terdegenerasi, dan syarat ini dapat diperiksa sesudah perubahan basis ke sembarang ekstensi medan. Terapkan Korolari \ref{prop:separable-discriminant} dan \ref{prop:etale-alg-characterization}, serta perhatikan bahwa jika $x \in A$ nilpoten, maka $\Tr_{A|F}(x)=0$.
	\end{hint}
\end{Exercises}
